% !TEX program = xelatex
% Thunder 四足机器人 Python SDK 开发手册
% 上海穹沛科技有限公司
% 2026年4月
\documentclass[a4paper,11pt,openany]{book}

% ============================================================
%  基础包
% ============================================================
\usepackage[UTF8, heading=true]{ctex}
\usepackage[margin=2.5cm, top=3cm, bottom=3cm]{geometry}
\usepackage{graphicx}
\usepackage{xcolor}
\usepackage{hyperref}
\usepackage{booktabs}
\usepackage{longtable}
\usepackage{array}
\usepackage{tabularx}
\usepackage{fancyhdr}
\usepackage{titlesec}
\usepackage{enumitem}
\usepackage{amsmath}
\usepackage{amssymb}
\usepackage{listings}
\usepackage{float}
\usepackage{pifont}
\usepackage{tocloft}

% ============================================================
%  颜色定义
% ============================================================
\definecolor{thunderblue}{RGB}{0, 82, 155}
\definecolor{thundergray}{RGB}{100, 100, 100}
\definecolor{codebg}{RGB}{245, 245, 245}
\definecolor{codeframe}{RGB}{200, 200, 200}
\definecolor{keywordcolor}{RGB}{0, 0, 180}
\definecolor{stringcolor}{RGB}{0, 128, 0}
\definecolor{commentcolor}{RGB}{128, 128, 128}
\definecolor{warnorange}{RGB}{200, 120, 0}

% ============================================================
%  超链接配置
% ============================================================
\hypersetup{
  colorlinks=true,
  linkcolor=thunderblue,
  urlcolor=thunderblue,
  citecolor=thunderblue,
  bookmarks=true,
  bookmarksnumbered=true,
  pdfauthor={上海穹沛科技有限公司},
  pdftitle={Thunder 四足机器人 Python SDK 开发手册},
  pdfsubject={Thunder SDK v0.1.0},
}

% ============================================================
%  代码高亮配置
% ============================================================
\lstset{
  basicstyle=\small\ttfamily,
  backgroundcolor=\color{codebg},
  frame=single,
  rulecolor=\color{codeframe},
  breaklines=true,
  breakatwhitespace=false,
  tabsize=4,
  showstringspaces=false,
  numbers=left,
  numberstyle=\tiny\color{thundergray},
  numbersep=8pt,
  xleftmargin=1.5em,
  framexleftmargin=1.5em,
  aboveskip=1em,
  belowskip=0.8em,
  literate={-}{-}1,
}

\lstdefinestyle{python}{
  language=Python,
  keywordstyle=\color{keywordcolor}\bfseries,
  stringstyle=\color{stringcolor},
  commentstyle=\color{commentcolor}\itshape,
  morekeywords={self, True, False, None, as, with, yield, async, await, f},
}

\lstdefinestyle{bash}{
  language=bash,
  keywordstyle=\color{keywordcolor}\bfseries,
  stringstyle=\color{stringcolor},
  commentstyle=\color{commentcolor}\itshape,
  morekeywords={ssh, ping, nmap, systemctl, pip, python, nc, scp, grep, cat, tail, ls, cd, export},
}

\lstdefinestyle{plain}{
  language={},
  keywordstyle={},
  stringstyle={},
  commentstyle={},
  numbers=none,
}

% ============================================================
%  页眉页脚
% ============================================================
\pagestyle{fancy}
\fancyhf{}
\fancyhead[LE]{\small\leftmark}
\fancyhead[RO]{\small\rightmark}
\fancyfoot[C]{\small\thepage}
\fancyfoot[LE,RO]{\small Thunder SDK v0.1.0}
\renewcommand{\headrulewidth}{0.4pt}
\renewcommand{\footrulewidth}{0.2pt}

% ============================================================
%  章节标题样式
% ============================================================
\titleformat{\chapter}[display]
  {\normalfont\huge\bfseries\color{thunderblue}}
  {\chaptertitlename\ \thechapter}{20pt}{\Huge}
\titlespacing*{\chapter}{0pt}{-20pt}{30pt}

% ============================================================
%  自定义命令
% ============================================================
\newcommand{\warning}[1]{%
  \begin{center}
  \fcolorbox{warnorange}{warnorange!10}{%
    \parbox{0.9\textwidth}{\textcolor{warnorange}{\textbf{注意：}}#1}%
  }
  \end{center}
}

\newcommand{\apimethod}[1]{%
  \paragraph{\texttt{#1}}%
}

% longtable 支持 booktabs
\setlength{\LTpre}{0.8em}
\setlength{\LTpost}{0.8em}

% ============================================================
%  文档开始
% ============================================================
\begin{document}

% ============================================================
%  封面
% ============================================================
\begin{titlepage}
\begin{center}

\vspace*{3cm}

{\Huge\bfseries\color{thunderblue} Thunder 四足机器人\\[0.5em] Python SDK 开发手册}

\vspace{2cm}

\includegraphics[width=3cm]{example-image-duck}% 占位图，可替换为公司 Logo

\vspace{2cm}

{\LARGE 上海穹沛科技有限公司}

\vspace{1cm}

{\Large\color{thundergray} Brainstem SDK for Python}

\vspace{2cm}

\begin{tabular}{rl}
\textbf{版本} & v0.1.0 \\
\textbf{日期} & 2026年4月 \\
\textbf{密级} & 内部文档 \\
\end{tabular}

\vfill

{\small\color{thundergray} Copyright \copyright\ 2026 上海穹沛科技有限公司. All rights reserved.}

\end{center}
\end{titlepage}

% ============================================================
%  版权页
% ============================================================
\thispagestyle{empty}
\vspace*{\fill}
\noindent\textbf{文档信息}\\[0.5em]
\begin{tabular}{ll}
\toprule
项目 & 说明 \\
\midrule
文档名称 & Thunder 四足机器人 Python SDK 开发手册 \\
版本号   & v0.1.0 \\
编写日期 & 2026年4月 \\
适用产品 & Thunder 四足轮腿机器人 \\
适用 SDK & brainstem\_sdk (Python) \\
\bottomrule
\end{tabular}

\vspace{2em}
\noindent\textbf{版本历史}\\[0.5em]
\begin{tabular}{llll}
\toprule
版本 & 日期 & 作者 & 说明 \\
\midrule
v0.1.0 & 2026-04 & 穹沛科技 & 初始版本 \\
\bottomrule
\end{tabular}

\vspace{2em}
\noindent{\small 本文档为上海穹沛科技有限公司内部技术文档，未经许可不得外传。}
\newpage

% ============================================================
%  目录
% ============================================================
\tableofcontents
\newpage

% ============================================================
%  概述 (来自 README.md)
% ============================================================
\chapter{概述}
\label{ch:overview}

Thunder 四足机器人 Python 开发套件，3 行代码控制机器人。

\section{安装}

\begin{lstlisting}[style=bash]
cd sdk/python
pip install -e .
\end{lstlisting}

\section{快速开始}

\begin{lstlisting}[style=python]
from brainstem_sdk import ThunderClient

# 连接机器人 (默认 192.168.66.190:13145)
dog = ThunderClient("192.168.66.190")

# 使能电机 -> 站起 -> 走路 -> 坐下 -> 禁用
dog.enable()
dog.stand_up()
dog.walk(vx=0.3)       # 30% 速度前进
dog.walk(vx=0, vyaw=0.2)  # 原地左转
dog.sit_down()
dog.disable()
\end{lstlisting}

\section{API 参考}

\subsection{连接}

\begin{lstlisting}[style=python]
# 默认连接
dog = ThunderClient("192.168.66.190")

# 自定义端口和超时
dog = ThunderClient("10.0.0.1", port=13145, timeout=5.0)

# 支持 with 语句
with ThunderClient("192.168.66.190") as dog:
    dog.stand_up()
\end{lstlisting}

\subsection{运动控制}

\begin{longtable}{lll}
\toprule
方法 & 说明 & 参数 \\
\midrule
\endhead
\texttt{dog.enable()} & 使能全部电机 & -- \\
\texttt{dog.disable()} & 禁用全部电机 & -- \\
\texttt{dog.stand\_up()} & 坐姿 $\to$ 站立 & -- \\
\texttt{dog.sit\_down()} & 站立 $\to$ 坐姿 & -- \\
\texttt{dog.walk(vx, vy, vyaw)} & 行走 & \parbox[t]{5cm}{vx: 前后 [-1,1]\\vy: 左右 [-1,1]\\vyaw: 旋转 [-1,1]} \\
\bottomrule
\end{longtable}

\noindent\textbf{walk() 参数约定：}
\begin{itemize}[nosep]
  \item \texttt{vx > 0} = 前进，\texttt{vx < 0} = 后退
  \item \texttt{vy > 0} = 向左，\texttt{vy < 0} = 向右
  \item \texttt{vyaw > 0} = 逆时针，\texttt{vyaw < 0} = 顺时针
  \item \texttt{dog.walk()} (无参数) = 原地停步
\end{itemize}

\subsection{状态查询}

\begin{lstlisting}[style=python]
# 获取当前状态
state = dog.get_state()
# 返回: "zero" | "grounded" | "standing" | "walking" | "transitioning"

# 获取当前策略
profile = dog.get_profile()
print(profile.current)     # "default"
print(profile.available)   # ["default", "mini"]

# 切换策略 (机器人必须在 grounded 状态)
dog.switch_profile("mini")
\end{lstlisting}

\subsection{诊断}

\begin{lstlisting}[style=python]
# 电压
voltages = dog.get_voltage()  # [42.1, 42.0, ...] 16个电机

# 电机状态
motors = dog.get_motor_status()
for m in motors:
    print(f"Joint {m.id}: online={m.online} temp={m.temperature}C")

# 清除电机故障
dog.clear_motor_fault()          # 清除全部
dog.clear_motor_fault([0, 1, 2]) # 只清前右腿

# 标零 (grounded 状态下)
dog.set_zero()
\end{lstlisting}

\subsection{实时数据流}

\begin{lstlisting}[style=python]
# IMU 数据 (~50Hz)
for imu in dog.listen_imu():
    print(f"gyro: {imu.gyroscope.x:.3f}, {imu.gyroscope.y:.3f}, {imu.gyroscope.z:.3f}")
    print(f"quat: w={imu.quaternion.w:.3f}")

# 关节数据
for joint in dog.listen_joint():
    print(f"Joint {joint.id}: pos={joint.position:.3f} rad")

# 状态变化
for state in dog.listen_state():
    print(f"State changed: {state}")
\end{lstlisting}

\section{关节编号}

\begin{verbatim}
前右 (FR)           前左 (FL)
  0 Hip               4 Hip
  1 Thigh             5 Thigh
  2 Calf              6 Calf
  3 Foot              7 Foot

后右 (RR)           后左 (RL)
  8 Hip              12 Hip
  9 Thigh            13 Thigh
 10 Calf             14 Calf
 11 Foot             15 Foot
\end{verbatim}

\section{状态机}

\begin{verbatim}
           Init
            |
            v
     +-- Grounded --+
     |      ^       |
     |      |       v
     |   SitDown  StandUp
     |      ^       |
     |      |       v
     |   Standing --+
     |      |
     |      v
     +-- Walking
\end{verbatim}

\begin{itemize}[nosep]
  \item \textbf{Grounded}：坐姿，安全状态。上电后自动进入。
  \item \textbf{Standing}：站立，可以接收 walk 指令。
  \item \textbf{Walking}：行走中。发送 \texttt{walk(0,0,0)} 回到 Standing。
\end{itemize}

\section{示例}

\begin{longtable}{ll}
\toprule
示例 & 说明 \\
\midrule
\endhead
\texttt{examples/01\_hello\_walk.py} & 最简走路示例 \\
\texttt{examples/02\_read\_imu.py} & 读取 IMU 数据流 \\
\texttt{examples/03\_motor\_diagnostics.py} & 电机诊断 \\
\texttt{examples/04\_keyboard\_control.py} & WASD 键盘控制 \\
\bottomrule
\end{longtable}

\section{网络要求}

\begin{itemize}[nosep]
  \item 开发机和机器人在同一局域网
  \item 默认 gRPC 端口：\textbf{13145}
  \item 无需安装额外中间件（纯 gRPC over HTTP/2）
\end{itemize}



% ============================================================
%  第 1 章：快速开始
% ============================================================
\chapter{快速开始}
\label{ch:quickstart}

本文档帮助您在 5 分钟内完成 Thunder SDK 的安装、连接和第一次控制。

\section{环境要求}

\begin{longtable}{ll}
\toprule
项目 & 要求 \\
\midrule
\endhead
Python & >= 3.8 \\
操作系统 & Windows / Linux / macOS \\
网络 & 开发机与机器人在同一局域网 \\
端口 & TCP 13145（gRPC over HTTP/2） \\
\bottomrule
\end{longtable}

\warning{不需要安装 ROS、消息中间件或其他额外组件。SDK 通过纯 gRPC 直连机器人。}

\section{安装}

\subsection{方式一：本地开发安装（推荐）}

\begin{lstlisting}[style=bash]
cd sdk/python
pip install -e .
\end{lstlisting}

\subsection{方式二：直接安装依赖}

\begin{lstlisting}[style=bash]
pip install grpcio>=1.50.0 protobuf>=4.21.0
\end{lstlisting}

安装完成后验证：

\begin{lstlisting}[style=bash]
python -c "from brainstem_sdk import ThunderClient; print('OK')"
\end{lstlisting}

输出 \texttt{OK} 即安装成功。

\section{三行代码控制机器人}

\begin{lstlisting}[style=python]
from brainstem_sdk import ThunderClient

dog = ThunderClient("192.168.66.190")
dog.walk(vx=0.3)  # 30% 速度前进
\end{lstlisting}

这三行代码做了以下事情：
\begin{enumerate}[nosep]
  \item 导入 SDK
  \item 通过 gRPC 连接到机器人（默认端口 13145）
  \item 发送前进指令
\end{enumerate}

\warning{实际使用中，您需要先 \texttt{enable()} 使能电机并 \texttt{stand\_up()} 站起来，才能执行 \texttt{walk()}。完整的操作流程见下方``第一次完整操作''。}

\section{首次连接检查清单}

在第一次连接机器人之前，请逐项确认：

\begin{itemize}[nosep]
  \item[$\square$] 机器人已上电，电池电压 > 42V
  \item[$\square$] 机器人主控板（S100P）已启动完成（约 30 秒）
  \item[$\square$] brainstem 服务已运行（\texttt{systemctl status brainstem}）
  \item[$\square$] 开发机与机器人在同一局域网（可以 \texttt{ping 192.168.66.190}）
  \item[$\square$] 遥控器已关闭或置于空闲状态（遥控器优先级高于 SDK）
  \item[$\square$] 机器人放置在平坦地面上，四足触地
\end{itemize}

\section{网络配置}

\begin{verbatim}
开发机 (PC/笔记本)                   Thunder 机器人
  192.168.66.xxx          <------>   192.168.66.190
       |                                  |
       +---- 同一局域网 (以太网/WiFi) -----+
                    端口: 13145/TCP
\end{verbatim}

\begin{longtable}{lll}
\toprule
参数 & 默认值 & 说明 \\
\midrule
\endhead
机器人 IP & \texttt{192.168.66.190} & S100P 主控板默认地址 \\
gRPC 端口 & \texttt{13145} & 无需修改 \\
协议 & gRPC (HTTP/2) & 无加密（局域网内使用） \\
超时 & 5.0 秒 & 可通过 \texttt{timeout} 参数调整 \\
\bottomrule
\end{longtable}

如果机器人 IP 不同，连接时指定即可：

\begin{lstlisting}[style=python]
dog = ThunderClient("10.0.0.100", port=13145, timeout=5.0)
\end{lstlisting}

\section{验证连接脚本}

将以下脚本保存为 \texttt{verify\_connection.py} 并运行：

\begin{lstlisting}[style=python]
"""验证与 Thunder 机器人的连接。"""

from brainstem_sdk import ThunderClient

ROBOT_IP = "192.168.66.190"

print(f"正在连接 Thunder @ {ROBOT_IP}:13145 ...")

try:
    dog = ThunderClient(ROBOT_IP, timeout=3.0)

    # 1. 查询状态机
    state = dog.get_state()
    print(f"[OK] 状态机: {state}")

    # 2. 查询电压
    voltages = dog.get_voltage()
    if voltages:
        avg_v = sum(voltages) / len(voltages)
        print(f"[OK] 平均电压: {avg_v:.1f}V ({len(voltages)} 个电机)")
    else:
        print("[OK] 电压查询成功（无数据）")

    # 3. 查询电机状态
    motors = dog.get_motor_status()
    online = sum(1 for m in motors if m.online)
    print(f"[OK] 电机在线: {online}/{len(motors)}")

    # 4. 查询当前策略
    profile = dog.get_profile()
    print(f"[OK] 当前策略: {profile.current}")
    print(f"     可用策略: {profile.available}")

    print("\n--- 连接验证通过 ---")

    dog.close()

except Exception as e:
    print(f"\n[FAIL] 连接失败: {e}")
    print("\n排查步骤:")
    print("  1. 确认机器人已上电并启动完成")
    print("  2. 确认网络连通: ping 192.168.66.190")
    print("  3. 确认服务运行: ssh sunrise@192.168.66.190 'systemctl status brainstem'")
    print("  4. 确认端口可达: 无防火墙阻断 13145/TCP")
\end{lstlisting}

运行：

\begin{lstlisting}[style=bash]
python verify_connection.py
\end{lstlisting}

预期输出：

\begin{lstlisting}[style=plain,numbers=none]
正在连接 Thunder @ 192.168.66.190:13145 ...
[OK] 状态机: grounded
[OK] 平均电压: 48.2V (16 个电机)
[OK] 电机在线: 16/16
[OK] 当前策略: default
     可用策略: ['default', 'mini']

--- 连接验证通过 ---
\end{lstlisting}

\section{第一次完整操作}

确认连接验证通过后，尝试完整的站起-行走-坐下流程：

\begin{lstlisting}[style=python]
import time
from brainstem_sdk import ThunderClient

dog = ThunderClient("192.168.66.190")

# 使能电机
dog.enable()
print(f"状态: {dog.get_state()}")   # grounded

# 站起来（约 3 秒过渡）
dog.stand_up()
time.sleep(4)
print(f"状态: {dog.get_state()}")   # standing

# 前进
dog.walk(vx=0.3)
time.sleep(3)

# 停步
dog.walk()    # 无参数 = 停步
time.sleep(1)

# 坐下（约 3 秒过渡）
dog.sit_down()
time.sleep(4)
print(f"状态: {dog.get_state()}")   # grounded

# 禁用电机
dog.disable()
dog.close()
\end{lstlisting}

\warning{首次操作建议在台架上进行，或有人在旁边准备随时用遥控器接管。遥控器具有最高优先级，任何时候拨动遥控器都可以立即接管控制权。}

\section{使用 with 语句}

SDK 支持上下文管理器，退出时自动关闭连接：

\begin{lstlisting}[style=python]
from brainstem_sdk import ThunderClient

with ThunderClient("192.168.66.190") as dog:
    state = dog.get_state()
    print(f"当前状态: {state}")
# 退出 with 块时自动调用 dog.close()
\end{lstlisting}

\section{常见问题}

\subsection{连接超时}

\begin{lstlisting}[style=plain,numbers=none]
grpc._channel._InactiveRpcError: ... UNAVAILABLE
\end{lstlisting}

\noindent\textbf{排查步骤}：
\begin{enumerate}[nosep]
  \item 确认机器人已上电且启动完成
  \item \texttt{ping 192.168.66.190} 确认网络连通
  \item SSH 登录确认服务运行：\texttt{ssh sunrise@192.168.66.190 'systemctl status brainstem'}
  \item 检查防火墙是否阻断了 13145/TCP
\end{enumerate}

\subsection{walk() 没反应}

必须按照状态机顺序操作：\texttt{enable()} $\to$ \texttt{stand\_up()} $\to$ 等待站稳 $\to$ \texttt{walk()}。

\begin{lstlisting}[style=python]
state = dog.get_state()
print(state)  # 如果不是 "standing"，walk() 不会生效
\end{lstlisting}

\subsection{电机不动}

\begin{enumerate}[nosep]
  \item 确认已调用 \texttt{dog.enable()}
  \item 遥控器在线时会抢占控制权------关闭或放下遥控器
  \item 查看电机状态：\texttt{dog.get\_motor\_status()}
\end{enumerate}

\subsection{操作被拒绝}

\begin{lstlisting}[style=plain,numbers=none]
grpc._channel._InactiveRpcError: ... transition in progress
\end{lstlisting}

状态机正在过渡中（如站起/坐下动画），请等待过渡完成后再发送新指令。



% ============================================================
%  第 2 章：硬件概览
% ============================================================
\chapter{硬件概览}
\label{ch:hardware}

本文档介绍 Thunder 四足轮腿机器人的硬件组成、传感器、执行器和通信架构。

\section{外观与尺寸}

\begin{longtable}{ll}
\toprule
参数 & 值 \\
\midrule
\endhead
类型 & 四足轮腿机器人 \\
自由度 & 16 DOF（4腿 $\times$ 3关节 + 4脚轮） \\
重量 & （待确认） \\
尺寸 (长$\times$宽$\times$高) & （待确认） \\
站立高度 & （待确认） \\
\bottomrule
\end{longtable}

\section{整体架构}

\begin{verbatim}
                    +--------------------------+
                    |       S100P 主控板        |
                    |    (RDK X5 / aarch64)    |
                    |     Ubuntu 22.04         |
                    +--+-----+-----+-----+----+
                       |     |     |     |
              +--------+  +--+--+  |  +--+--------+
              |           |     |  |  |            |
         +----+----+  +---+--+ +--+--+--+  +------+------+
         | HI91 IMU|  | PCAN | | 以太网  |  | SBUS/USB   |
         | (50Hz)  |  | CAN  | | gRPC   |  | 遥控器      |
         +---------+  +---+--+ +---------+  +-------------+
                          |
            +-------------+-------------+
            |             |             |
      +-----+-----+ +----+----+ +------+------+
      | Robstride  | | Robstride| | Robstride  |
      | 电机 x4    | | 电机 x4  | | 电机 x4+4  |
      | (FR leg)   | | (FL leg) | | (RR+RL)    |
      +------------+ +----------+ +-------------+
\end{verbatim}

\section{计算平台}

\begin{longtable}{ll}
\toprule
参数 & 规格 \\
\midrule
\endhead
主控板 & 地瓜机器人 S100P (RDK X5) \\
处理器 & Horizon J6, aarch64 \\
AI 算力 & Nash BPU, 128 TOPS \\
操作系统 & Ubuntu 22.04 \\
中间件 & ROS2 Humble (可选) \\
控制程序 & brainstem (Dart runtime) \\
\bottomrule
\end{longtable}

\noindent 主控板负责：
\begin{itemize}[nosep]
  \item 50Hz 实时控制循环
  \item ONNX 神经网络推理（步态策略）
  \item gRPC 服务端（接收 SDK 指令）
  \item IMU 数据采集与融合
  \item CAN 总线电机通信
  \item 低压保护与安全监控
\end{itemize}

\section{关节与执行器}

\subsection{关节布局}

Thunder 共有 \textbf{16 个自由度}，分为 4 条腿，每条腿包含 3 个旋转关节和 1 个脚轮：

\begin{verbatim}
        前 (Front)
   FL +---------+ FR
      |         |
      |  body   |
      |         |
   RL +---------+ RR
        后 (Rear)

每条腿的关节（由近到远）:
  Hip   -- 髋关节（横滚方向，内外摆动）
  Thigh -- 大腿关节（俯仰方向，前后摆动）
  Calf  -- 小腿关节（俯仰方向，前后摆动）
  Foot  -- 脚轮（连续旋转，轮式驱动）
\end{verbatim}

\subsection{电机规格}

\begin{longtable}{lll}
\toprule
参数 & 腿部关节 (Hip/Thigh/Calf) & 脚轮 (Foot) \\
\midrule
\endhead
电机 & Robstride & Robstride \\
总线 & PCAN CAN & PCAN CAN \\
最大力矩 & 50 Nm & 30 Nm \\
控制模式 & PD 位置控制 & 速度控制 (Kp=0) \\
反馈 & 位置 / 速度 / 力矩 / 温度 & 位置 / 速度 / 力矩 / 温度 \\
\bottomrule
\end{longtable}

\subsection{PD 控制参数（default 策略）}

行走推理阶段的 PD 增益：

\begin{longtable}{llll}
\toprule
关节类型 & Kp & Kd & 动作缩放 \\
\midrule
\endhead
Hip (髋) & 80 & 10 & 0.15 \\
Thigh (大腿) & 100 & 10 & 0.25 \\
Calf (小腿) & 120 & 15 & 0.25 \\
Foot (脚轮) & 0 & 1 & 5.0 \\
\bottomrule
\end{longtable}

\warning{脚轮的 Kp=0 表示不做位置控制，仅通过速度指令驱动。}

站起/坐下阶段使用更高的刚度：Kp=200, Kd=8（全部 16 个关节统一）。

\section{传感器}

\subsection{惯性测量单元 (IMU)}

\begin{longtable}{ll}
\toprule
参数 & 规格 \\
\midrule
\endhead
型号 & HI91 \\
数据 & 三轴陀螺仪 + 四元数姿态 \\
频率 & \textasciitilde 50Hz \\
安装位置 & 机身中央 \\
坐标系 & x=前, y=左, z=上 \\
\bottomrule
\end{longtable}

IMU 提供两类数据：
\begin{itemize}[nosep]
  \item \textbf{陀螺仪} (\texttt{gyroscope})：三轴角速度 (rad/s)，用于步态平衡
  \item \textbf{四元数} (\texttt{quaternion})：当前姿态 (w,x,y,z Hamilton 约定)，用于重力投影
\end{itemize}

SDK 获取方式：

\begin{lstlisting}[style=python]
for imu in dog.listen_imu():
    print(f"陀螺仪: x={imu.gyroscope.x:.3f} y={imu.gyroscope.y:.3f} z={imu.gyroscope.z:.3f}")
    print(f"姿态:   w={imu.quaternion.w:.3f} x={imu.quaternion.x:.3f}")
\end{lstlisting}

\subsection{关节编码器}

每个电机内置编码器，提供：

\begin{longtable}{lll}
\toprule
数据 & 单位 & 说明 \\
\midrule
\endhead
位置 (position) & rad & 当前关节角度 \\
速度 (velocity) & rad/s & 当前关节角速度 \\
力矩 (torque) & Nm & 当前输出力矩 \\
\bottomrule
\end{longtable}

SDK 获取方式：

\begin{lstlisting}[style=python]
for joint in dog.listen_joint():
    print(f"关节{joint.id}: pos={joint.position:.3f}rad vel={joint.velocity:.3f}rad/s")
\end{lstlisting}

\subsection{总线电压}

通过 CAN 总线读取电机端电压，反映电池状态：

\begin{lstlisting}[style=python]
voltages = dog.get_voltage()  # 返回 16 个电机的电压值 (V)
avg = sum(voltages) / len(voltages)
print(f"平均电压: {avg:.1f}V")
\end{lstlisting}

\section{通信接口}

\subsection{接口一览}

\begin{longtable}{llll}
\toprule
接口 & 协议 & 用途 & 方向 \\
\midrule
\endhead
以太网 & gRPC (HTTP/2) & SDK 远程控制 & 双向 \\
PCAN CAN & CAN 2.0B & 电机指令/反馈 & 双向 \\
SBUS 串口 & SBUS & YUNZHUO 遥控器 & 输入 \\
USB & HID & Xbox 手柄 (备选) & 输入 \\
串口 & UART & HI91 IMU & 输入 \\
\bottomrule
\end{longtable}

\subsection{gRPC 服务}

\begin{longtable}{ll}
\toprule
参数 & 值 \\
\midrule
\endhead
端口 & 13145 \\
协议 & gRPC (Protobuf) \\
加密 & 无（局域网内使用） \\
服务名 & CMS (Central Motor System) \\
\bottomrule
\end{longtable}

gRPC 提供的 RPC 接口：

\begin{longtable}{lll}
\toprule
类别 & 方法 & 说明 \\
\midrule
\endhead
电机控制 & \texttt{Enable} / \texttt{Disable} & 使能/禁用电机 \\
运动指令 & \texttt{Walk} / \texttt{StandUp} / \texttt{SitDown} & 状态机动作 \\
状态查询 & \texttt{GetCmsState} / \texttt{GetVoltage} / \texttt{GetMotorStatus} & 读取当前状态 \\
策略管理 & \texttt{GetProfile} / \texttt{SwitchProfile} & 切换运动策略 \\
诊断 & \texttt{ClearMotorFault} / \texttt{SetZero} & 清错/标零 \\
数据流 & \texttt{ListenImu} / \texttt{ListenJoint} / \texttt{ListenCmsState} & 实时推送 \\
\bottomrule
\end{longtable}

\subsection{控制器输入}

Thunder 支持两种物理控制器，均通过 \texttt{ControlArbiter} 仲裁控制权：

\begin{longtable}{llll}
\toprule
控制器 & 接口 & 优先级 & 说明 \\
\midrule
\endhead
YUNZHUO 遥控器 & SBUS 串口 & 高 & 硬件遥控，最高优先级 \\
Xbox 手柄 & USB HID & 高 & Linux Joystick FFI 读取 \\
SDK (gRPC) & 以太网 & 低 & 被遥控器/手柄抢占 \\
\bottomrule
\end{longtable}

\warning{遥控器/手柄信号到达时立即抢占 SDK 控制权。遥控器松开 3 秒后自动释放控制权回 SDK。}

\section{电源与电池}

\begin{longtable}{ll}
\toprule
参数 & 值 \\
\midrule
\endhead
供电 & 锂电池组 \\
工作电压范围 & （待确认满电电压）\textasciitilde{} 42V \\
低压保护阈值 & \textbf{42V} \\
低压保护行为 & 自动坐下 (graceful sit down) \\
\bottomrule
\end{longtable}

低压保护流程：
\begin{enumerate}[nosep]
  \item 控制循环每帧检测所有电机总线电压
  \item 当最低电压 $<$ 42V 时触发保护
  \item 机器人自动执行坐下动作
  \item 等待完全坐稳（Grounded 状态）后才禁用电机
  \item 不会直接断电，确保安全着地
\end{enumerate}

\section{控制频率}

\begin{longtable}{lll}
\toprule
环节 & 频率 & 周期 \\
\midrule
\endhead
主控制循环 & 50 Hz & 20 ms \\
ONNX 推理 & 50 Hz & 20 ms（每帧一次） \\
IMU 采样 & \textasciitilde 50 Hz & \textasciitilde 20 ms \\
CAN 电机通信 & 50 Hz & 20 ms \\
gRPC 状态推送 & 50 Hz & 20 ms \\
\bottomrule
\end{longtable}

整个控制链路在 20ms 的时间窗口内完成：

\begin{verbatim}
IMU采样 -> 构建观测 -> ONNX推理 -> PD计算 -> CAN发送 -> 电机执行
                        |
                  全部在 20ms 内完成
\end{verbatim}



% ============================================================
%  第 3 章：坐标系与关节
% ============================================================
\chapter{坐标系与关节}
\label{ch:coordinates}

本文档定义 Thunder 机器人的坐标系约定、16 个关节的编号规则、运动范围和典型姿态。

\section{机体坐标系}

Thunder 使用右手坐标系，原点位于机身几何中心：

\begin{verbatim}
              +Z (上)
               |
               |
               |
               +-------> +X (前)
              /
             /
            +Y (左)
\end{verbatim}

\begin{longtable}{lll}
\toprule
轴 & 正方向 & 说明 \\
\midrule
\endhead
X & 前方 (Forward) & 机器人头部方向 \\
Y & 左方 (Left) & 机器人左侧方向 \\
Z & 上方 (Up) & 垂直向上 \\
\bottomrule
\end{longtable}

\subsection{正面视图}

\begin{verbatim}
               +Z
                |
        FL      |      FR
    +---+-------+-------+---+
    |   |       |       |   |
    |   |     body      |   |
    |   |       |       |   |
    +---+-------+-------+---+
        RL      |      RR
                |
         +Y <--+
\end{verbatim}

\subsection{侧面视图（右侧）}

\begin{verbatim}
         +Z
          |           前进方向 --->
          |
          +------ body ------+
          |                  |
         Hip               Hip
          |                  |
        Thigh             Thigh
          |                  |
        Calf              Calf
          |                  |
        [Foot]           [Foot]
    ----地面-----------------------
          |
          +-------> +X (前)

    后腿 (RR)           前腿 (FR)
\end{verbatim}

\subsection{walk() 指令与坐标系的关系}

\begin{longtable}{llll}
\toprule
参数 & 正值 & 负值 & 对应轴 \\
\midrule
\endhead
\texttt{vx} & 前进 & 后退 & +X / -X \\
\texttt{vy} & 向左 & 向右 & +Y / -Y \\
\texttt{vyaw} & 逆时针 & 顺时针 & 绕 +Z 轴 \\
\bottomrule
\end{longtable}

\section{关节编号总表}

Thunder 共 16 个自由度，按腿分组编号：

{\small
\begin{longtable}{clllllll}
\toprule
索引 & 名称 & 腿 & 关节类型 & URDF 名称 & 范围 (rad) & 正方向 & 最大力矩 \\
\midrule
\endhead
0 & FR Hip & 前右 & 髋 (横滚) & \texttt{fr\_hip\_joint} & -0.41 \textasciitilde{} +0.41 & 外展 & 50 Nm \\
1 & FR Thigh & 前右 & 大腿 (俯仰) & \texttt{fr\_thigh\_joint} & -3.14 \textasciitilde{} 0 & 后摆 & 50 Nm \\
2 & FR Calf & 前右 & 小腿 (俯仰) & \texttt{fr\_calf\_joint} & 0 \textasciitilde{} +3.14 & 伸展 & 50 Nm \\
3 & FR Foot & 前右 & 脚轮 (连续) & \texttt{fr\_foot\_joint} & 连续旋转 & -- & 30 Nm \\
4 & FL Hip & 前左 & 髋 (横滚) & \texttt{fl\_hip\_joint} & -0.41 \textasciitilde{} +0.41 & 外展 & 50 Nm \\
5 & FL Thigh & 前左 & 大腿 (俯仰) & \texttt{fl\_thigh\_joint} & 0 \textasciitilde{} +3.14 & 前摆 & 50 Nm \\
6 & FL Calf & 前左 & 小腿 (俯仰) & \texttt{fl\_calf\_joint} & -3.14 \textasciitilde{} 0 & 收缩 & 50 Nm \\
7 & FL Foot & 前左 & 脚轮 (连续) & \texttt{fl\_foot\_joint} & 连续旋转 & -- & 30 Nm \\
8 & RR Hip & 后右 & 髋 (横滚) & \texttt{rr\_hip\_joint} & -0.41 \textasciitilde{} +0.41 & 外展 & 50 Nm \\
9 & RR Thigh & 后右 & 大腿 (俯仰) & \texttt{rr\_thigh\_joint} & 0 \textasciitilde{} +3.14 & 前摆 & 50 Nm \\
10 & RR Calf & 后右 & 小腿 (俯仰) & \texttt{rr\_calf\_joint} & -6.28 \textasciitilde{} 0 & 收缩 & 50 Nm \\
11 & RR Foot & 后右 & 脚轮 (连续) & \texttt{rr\_foot\_joint} & 连续旋转 & -- & 30 Nm \\
12 & RL Hip & 后左 & 髋 (横滚) & \texttt{rl\_hip\_joint} & -0.41 \textasciitilde{} +0.41 & 外展 & 50 Nm \\
13 & RL Thigh & 后左 & 大腿 (俯仰) & \texttt{rl\_thigh\_joint} & -3.14 \textasciitilde{} 0 & 后摆 & 50 Nm \\
14 & RL Calf & 后左 & 小腿 (俯仰) & \texttt{rl\_calf\_joint} & 0 \textasciitilde{} +6.28 & 伸展 & 50 Nm \\
15 & RL Foot & 后左 & 脚轮 (连续) & \texttt{rl\_foot\_joint} & 连续旋转 & -- & 30 Nm \\
\bottomrule
\end{longtable}
}

\section{关节分组}

\subsection{按腿分组}

\begin{verbatim}
前右 (FR): [0, 1, 2, 3]    前左 (FL): [4, 5, 6, 7]
后右 (RR): [8, 9, 10, 11]  后左 (RL): [12, 13, 14, 15]
\end{verbatim}

\subsection{按关节类型分组}

\begin{verbatim}
所有 Hip:   [0, 4, 8, 12]   -- 髋关节，控制腿的内外摆动
所有 Thigh: [1, 5, 9, 13]   -- 大腿关节，控制腿的前后摆动
所有 Calf:  [2, 6, 10, 14]  -- 小腿关节，控制膝盖弯曲
所有 Foot:  [3, 7, 11, 15]  -- 脚轮，控制轮式移动
\end{verbatim}

\subsection{俯瞰图 (Top View)}

\begin{verbatim}
        前 (Head)
          X+
          |
   FL[4-7]|  FR[0-3]
     +-----+-----+
     |     |     |
     |   body    |
 Y+--+     O     +-- Y-
     |           |
     |           |
     +-----+-----+
   RL[12-15] RR[8-11]
          |
          X-
        后 (Tail)
\end{verbatim}

\subsection{侧面图------单腿结构}

\begin{verbatim}
        body
         |
    [0] Hip -------- 横滚轴 (绕 X 轴)
         |
         |  上臂
         |
    [1] Thigh ------ 俯仰轴 (绕 Y 轴)
         |
         |  下臂
         |
    [2] Calf ------- 俯仰轴 (绕 Y 轴)
         |
         |  小腿
         |
    [3] Foot ------- 连续旋转 (轮轴)
        (o)          脚轮触地
    ============== 地面
\end{verbatim}

\section{典型姿态}

\subsection{站立姿态 (Standing Pose)}

站立姿态是机器人正常直立时各关节的目标角度：

\begin{lstlisting}[style=python]
standing_pose = [
    -0.1, -0.8,  1.8,  0,    # FR: hip, thigh, calf, foot
     0.1,  0.8, -1.8,  0,    # FL: hip, thigh, calf, foot
     0.1,  0.8, -1.8,  0,    # RR: hip, thigh, calf, foot
    -0.1, -0.8,  1.8,  0,    # RL: hip, thigh, calf, foot
]
\end{lstlisting}

各关节含义：

\begin{longtable}{lll}
\toprule
关节 & FR 值 & 含义 \\
\midrule
\endhead
Hip & -0.1 rad ($-5.7\degree$) & 微微内收，保持稳定 \\
Thigh & -0.8 rad ($-45.8\degree$) & 大腿向后摆约 46 度 \\
Calf & 1.8 rad ($103.1\degree$) & 小腿弯曲约 103 度 \\
Foot & 0 rad & 脚轮不转 \\
\bottomrule
\end{longtable}

\warning{注意左右对称：FR 和 RL 的符号相同，FL 和 RR 的符号相同（镜像关系）。}

\subsection{坐下姿态 (Sitting Pose)}

坐下姿态即电机零位，所有关节角度为 0：

\begin{lstlisting}[style=python]
sitting_pose = [0, 0, 0, 0,   # FR
                0, 0, 0, 0,   # FL
                0, 0, 0, 0,   # RR
                0, 0, 0, 0]   # RL
\end{lstlisting}

这是机器人趴在地面上的姿态，也是电机标零后的参考位。

\subsection{姿态可视化}

\begin{verbatim}
站立 (Standing)                  坐下 (Sitting / Zero)

    +-------+                        +-------+
    |  body |                        |  body |
    +--+-+--+                        +--+-+--+
       | |                             /   \
      /   \                           /     \
     /     \                      ---/-------\---
    /       \                        (地面)
   |         |
   |         |
  (o)       (o)
---地面-----------
\end{verbatim}

\section{PD 控制参数}

每个关节的运动由 PD 控制器驱动：

\begin{verbatim}
torque = Kp * (target_pos - current_pos) + Kd * (target_vel - current_vel)
\end{verbatim}

\subsection{推理阶段（行走时）}

\begin{longtable}{lllll}
\toprule
索引 & 关节 & Kp & Kd & actionScale \\
\midrule
\endhead
0, 4, 8, 12 & Hip & 80 & 10 & 0.15 \\
1, 5, 9, 13 & Thigh & 100 & 10 & 0.25 \\
2, 6, 10, 14 & Calf & 120 & 15 & 0.25 \\
3, 7, 11, 15 & Foot & 0 & 1 & 5.0 \\
\bottomrule
\end{longtable}

\texttt{actionScale} 是 RL 策略输出到实际弧度的缩放系数。策略输出值经过以下变换得到目标关节角度：

\begin{verbatim}
target_position = action_output * actionScale + standingPose
\end{verbatim}

\subsection{过渡阶段（站起/坐下）}

站起和坐下使用统一的高刚度参数，确保平稳过渡：

\begin{longtable}{lll}
\toprule
参数 & 值 & 说明 \\
\midrule
\endhead
Kp & 200 & 全部 16 个关节 \\
Kd & 8 & 全部 16 个关节 \\
\bottomrule
\end{longtable}

\section{关节速度缩放}

IMU 陀螺仪和关节速度在输入观测向量前会进行缩放：

\begin{longtable}{lll}
\toprule
参数 & 缩放系数 & 说明 \\
\midrule
\endhead
\texttt{imuGyroscopeScale} & 0.25 & IMU 陀螺仪数据缩放 \\
\texttt{jointVelocityScale} & {[}0.05, 0.05, 0.05, 0.05{]} & 各类型关节速度缩放 \\
\bottomrule
\end{longtable}

\section{编程中的关节操作}

\subsection{读取关节数据}

\begin{lstlisting}[style=python]
# 实时流式读取
for joint in dog.listen_joint():
    print(f"关节 {joint.id}: "
          f"位置={joint.position:.3f}rad "
          f"速度={joint.velocity:.3f}rad/s "
          f"力矩={joint.torque:.2f}Nm")
\end{lstlisting}

\subsection{诊断单条腿}

\begin{lstlisting}[style=python]
# 检查前右腿 (FR) 的 4 个电机
motors = dog.get_motor_status()
fr_motors = [m for m in motors if m.id in [0, 1, 2, 3]]
for m in fr_motors:
    print(f"  关节{m.id}: online={m.online} temp={m.temperature:.1f}C")
\end{lstlisting}

\subsection{清除单腿故障}

\begin{lstlisting}[style=python]
# 只清除前右腿的电机故障
dog.clear_motor_fault(joint_ids=[0, 1, 2, 3])
\end{lstlisting}



% ============================================================
%  第 4 章：状态机与运动控制
% ============================================================
\chapter{状态机与运动控制}
\label{ch:fsm}

本文档详细说明 Thunder 机器人的有限状态机 (FSM)、状态转换规则、运动控制接口和安全机制。

\section{状态机总览}

Thunder 的运动控制由一个严格的有限状态机管理。所有 SDK 指令必须遵循状态机规则，非法的状态转换会被拒绝。

\subsection{状态图}

\begin{verbatim}
                   +---------+
                   |  Zero   |     上电初始状态
                   +----+----+
                        |
                      Init()
                        |
                        v
                   +---------+
          +------->|Grounded |<--------+
          |        +----+----+         |
          |             |              |
       SitDown()    StandUp()     SitDown 完成
       完成             |          或 Fault
          |             v              |
          |     +-------+-------+      |
          |     | Transitioning |------+
          |     |  (StandUp)    |
          |     +-------+-------+
          |             |
          |          Done()
          |             |
          |             v
          |        +---------+
          +--------|Standing |<--------+
          |        +----+----+         |
          |             |              |
       SitDown()    Walk(dir)     Idle() / Fault
          |             |          或 StandUp()
          |             v              |
          |     +-------+-------+      |
          |     | Transitioning |      |
          |     |  (SitDown)    |      |
          |     +-------+-------+      |
          |             |              |
          |             v              |
          |        +---------+         |
          +--------|Walking  |---------+
                   +---------+
\end{verbatim}

\subsection{五个状态}

\begin{longtable}{lll}
\toprule
状态 & SDK 返回值 & 说明 \\
\midrule
\endhead
\textbf{Zero} & \texttt{"zero"} & 上电初始态，FSM 尚未初始化。不接受任何运动指令。 \\
\textbf{Grounded} & \texttt{"grounded"} & 坐姿/趴地。安全状态，电机可使能但不运动。 \\
\textbf{Standing} & \texttt{"standing"} & 直立站稳。可以接收 \texttt{walk()} 或 \texttt{sit\_down()} 指令。 \\
\textbf{Walking} & \texttt{"walking"} & 行走中。RL 策略以 50Hz 推理驱动步态。 \\
\textbf{Transitioning} & \texttt{"transitioning"} & 过渡中（站起/坐下/动作播放）。拒绝新的运动指令。 \\
\bottomrule
\end{longtable}

\section{状态转换详解}

\subsection{Zero $\to$ Grounded}

\begin{longtable}{ll}
\toprule
项目 & 说明 \\
\midrule
\endhead
触发 & 系统自动初始化 (Init) \\
耗时 & 即时 \\
用户操作 & 无需操作，上电后自动完成 \\
\bottomrule
\end{longtable}

上电启动 brainstem 服务后，FSM 自动从 Zero 进入 Grounded。

\subsection{Grounded $\to$ Standing（站起）}

\begin{longtable}{ll}
\toprule
项目 & 说明 \\
\midrule
\endhead
触发 & \texttt{dog.stand\_up()} \\
中间状态 & Transitioning (StandUp) \\
耗时 & 150 ticks $\times$ 20ms = \textbf{3 秒} \\
过程 & 从当前姿态线性插值到 standingPose \\
PD 参数 & Kp=200, Kd=8（高刚度，确保平稳） \\
\bottomrule
\end{longtable}

站起过程是一段 150 帧的线性插值动画：

\begin{verbatim}
第 0 帧:   关节角度 = 当前位置 (sittingPose)
第 75 帧:  关节角度 = 50% 位置 (中间态)
第 150 帧: 关节角度 = standingPose (完全站立)
\end{verbatim}

\subsection{Standing $\to$ Walking（开始行走）}

\begin{longtable}{ll}
\toprule
项目 & 说明 \\
\midrule
\endhead
触发 & \texttt{dog.walk(vx, vy, vyaw)} \\
耗时 & 即时切换 \\
过程 & RL 策略接管，以 50Hz 推理输出关节指令 \\
\bottomrule
\end{longtable}

\subsection{Walking 中更新方向}

\begin{longtable}{ll}
\toprule
项目 & 说明 \\
\midrule
\endhead
触发 & 再次调用 \texttt{dog.walk(vx, vy, vyaw)} \\
耗时 & 即时生效（下一帧） \\
过程 & 仅更新方向向量，不切换状态 \\
\bottomrule
\end{longtable}

\subsection{Walking $\to$ Standing（停步）}

\begin{longtable}{ll}
\toprule
项目 & 说明 \\
\midrule
\endhead
触发 & \texttt{dog.walk()} 无参数，或发送全零后超时自动触发 Idle \\
中间状态 & Transitioning (StandUp)，从当前行走姿态插值到 standingPose \\
耗时 & \textasciitilde 3 秒 \\
\bottomrule
\end{longtable}

\subsection{Standing $\to$ Grounded（坐下）}

\begin{longtable}{ll}
\toprule
项目 & 说明 \\
\midrule
\endhead
触发 & \texttt{dog.sit\_down()} \\
中间状态 & Transitioning (SitDown) \\
耗时 & 150 ticks $\times$ 20ms = \textbf{3 秒} \\
过程 & 从 standingPose 线性插值到 sittingPose (全零) \\
\bottomrule
\end{longtable}

\subsection{Walking $\to$ Grounded（行走中直接坐下）}

\begin{longtable}{ll}
\toprule
项目 & 说明 \\
\midrule
\endhead
触发 & 在 Walking 状态调用 \texttt{dog.sit\_down()} \\
过程 & 复合过渡：先自动站稳 (StandUp)，再坐下 (SitDown) \\
耗时 & \textasciitilde 6 秒（两段过渡串联） \\
\bottomrule
\end{longtable}

\section{非法操作处理}

在错误的状态下发送指令不会导致错误，但指令会被静默忽略：

\begin{longtable}{lll}
\toprule
当前状态 & 发送的指令 & 结果 \\
\midrule
\endhead
Grounded & \texttt{walk()} & 忽略。必须先 \texttt{stand\_up()} \\
Grounded & \texttt{sit\_down()} & 忽略。已在坐姿 \\
Standing & \texttt{stand\_up()} & 忽略。已站立 \\
Walking & \texttt{stand\_up()} & 触发从行走过渡到站立 \\
Transitioning & 任何运动指令 & \textbf{拒绝}。过渡期间不接受新指令 \\
Zero & 任何运动指令 & 忽略。FSM 未初始化 \\
\bottomrule
\end{longtable}

\noindent\textbf{最佳实践}：在发送指令前先查询状态，确认状态机处于正确状态。

\begin{lstlisting}[style=python]
state = dog.get_state()
if state == "standing":
    dog.walk(vx=0.3)
elif state == "grounded":
    dog.stand_up()
\end{lstlisting}

\section{walk() 速度约定}

\subsection{参数说明}

\begin{lstlisting}[style=python]
dog.walk(vx=0.0, vy=0.0, vyaw=0.0)
\end{lstlisting}

\begin{longtable}{lllll}
\toprule
参数 & 范围 & 正值 & 负值 & 零值 \\
\midrule
\endhead
\texttt{vx} & {[}-1.0, 1.0{]} & 前进 (+X) & 后退 (-X) & 不动 \\
\texttt{vy} & {[}-1.0, 1.0{]} & 向左 (+Y) & 向右 (-Y) & 不动 \\
\texttt{vyaw} & {[}-1.0, 1.0{]} & 逆时针 (+Z) & 顺时针 (-Z) & 不转 \\
\bottomrule
\end{longtable}

三个参数均为\textbf{归一化速度}，1.0 对应最大速度，具体物理速度由 RL 策略决定。

\subsection{常见用法}

\begin{lstlisting}[style=python]
# 前进
dog.walk(vx=0.5)

# 后退
dog.walk(vx=-0.3)

# 向左平移
dog.walk(vy=0.3)

# 原地逆时针转
dog.walk(vyaw=0.3)

# 边走边转（前进 + 左转）
dog.walk(vx=0.3, vyaw=0.2)

# 停步（发送零速度）
dog.walk()
\end{lstlisting}

\subsection{运动方向参考图}

\begin{verbatim}
                  vx > 0 (前进)
                     ^
                     |
                     |
    vy > 0 (左) <----O----> vy < 0 (右)
                     |
                     |
                     v
                  vx < 0 (后退)

            vyaw > 0        vyaw < 0
              (CCW)           (CW)
               <--             -->
              /   \           /   \
             |  O  |         |  O  |
              \   /           \   /
               -->             <--
\end{verbatim}

\section{电机使能/禁用生命周期}

\subsection{完整生命周期}

\begin{verbatim}
  上电 (Power On)
       |
       v
  [电机禁用] ---enable()---> [电机使能]
       ^                         |
       |                     正常使用:
   disable()               stand_up/walk/sit_down
       |                         |
       +-------disable()----<----+
       |
       v
  [电机禁用]
       |
  下电 (Power Off)
\end{verbatim}

\subsection{关键规则}

\begin{enumerate}[nosep]
  \item \textbf{上电后电机默认禁用}。必须调用 \texttt{enable()} 才能运动。
  \item \textbf{\texttt{disable()} 随时可调用}。会立即释放所有电机力矩（机器人会瘫软）。
  \item \textbf{先坐下再禁用}。建议流程：\texttt{sit\_down()} $\to$ 等待 Grounded $\to$ \texttt{disable()}。
  \item \textbf{低压自动禁用}。电压低于 42V 时，系统会先坐下再禁用，不会直接断电。
\end{enumerate}

\subsection{安全的关机流程}

\begin{lstlisting}[style=python]
import time

# 确保从任何状态安全关闭
state = dog.get_state()

if state == "walking":
    dog.sit_down()      # Walking -> StandUp -> SitDown (复合)
    time.sleep(7)       # 等待两段过渡完成
elif state == "standing":
    dog.sit_down()
    time.sleep(4)
elif state == "transitioning":
    time.sleep(4)       # 等待当前过渡完成

# 确认已坐下
assert dog.get_state() == "grounded"

# 安全禁用
dog.disable()
\end{lstlisting}

\section{控制权仲裁 (ControlArbiter)}

\subsection{优先级机制}

Thunder 支持多个控制源同时存在，通过 ControlArbiter 仲裁控制权：

\begin{verbatim}
优先级高                                     优先级低
+-----------+     +-----------+     +-----------+
| YUNZHUO   | >>> | Xbox 手柄 | >>> | SDK gRPC  |
| 遥控器    |     |           |     |           |
+-----------+     +-----------+     +-----------+
\end{verbatim}

\subsection{仲裁规则}

\begin{longtable}{ll}
\toprule
场景 & 行为 \\
\midrule
\endhead
SDK 控制中，遥控器摇杆动 & SDK 被立即抢占，遥控器接管 \\
遥控器控制中，SDK 发指令 & SDK 指令被忽略 \\
遥控器松手不动 & 3 秒超时后自动释放，SDK 恢复控制权 \\
SDK 控制中，遥控器关机 & SDK 保持控制权，不受影响 \\
\bottomrule
\end{longtable}

\subsection{超时参数}

\begin{longtable}{llll}
\toprule
参数 & 默认值 & 环境变量 & 说明 \\
\midrule
\endhead
\texttt{arbiterTimeout} & 3 秒 & \texttt{HAN\_DOG\_ARBITER\_TIMEOUT} & 遥控器释放控制权的超时时间 \\
\bottomrule
\end{longtable}

\subsection{SDK 开发建议}

\begin{enumerate}[nosep]
  \item \textbf{首次调试时关闭遥控器}，避免控制权冲突
  \item \textbf{保留遥控器作为安全手段}，紧急情况下随时接管
  \item \textbf{不要轮询状态来检测控制权}，SDK 指令被拒绝时不会报错，只是静默忽略
\end{enumerate}

\section{故障处理}

\subsection{Fault 状态转换}

当系统检测到异常时，会触发 Fault 动作。不同状态下的 Fault 行为：

\begin{longtable}{ll}
\toprule
当前状态 & Fault 行为 \\
\midrule
\endhead
Grounded & 忽略（已在安全状态） \\
Standing & 忽略（已在安全状态） \\
Walking & 自动过渡到 StandUp（先站稳） \\
Transitioning (StandUp) & 中止站起，转为 SitDown \\
Transitioning (SitDown) & 强制进入 Grounded（避免死循环） \\
\bottomrule
\end{longtable}

\subsection{电机故障清除}

\begin{lstlisting}[style=python]
# 查看电机状态
motors = dog.get_motor_status()
for m in motors:
    if m.errors:
        print(f"关节 {m.id} 故障: {m.errors}")

# 清除全部故障
dog.clear_motor_fault()

# 清除指定关节的故障
dog.clear_motor_fault(joint_ids=[0, 1, 2])
\end{lstlisting}

\subsection{低压保护}

当任一电机总线电压低于 42V 时触发：

\begin{verbatim}
检测低压 -> 触发 graceful sit_down -> 等待 Grounded -> 禁用电机
\end{verbatim}

低压保护是\textbf{不可逆的}（单次触发后标记 \texttt{lowVoltageTriggered}），需要重新上电才能恢复。

\section{策略切换}

Thunder 支持多种运动策略（Profile），每种策略包含不同的 ONNX 模型和 PD 参数。

\subsection{查询和切换}

\begin{lstlisting}[style=python]
# 查看当前策略
profile = dog.get_profile()
print(f"当前: {profile.current}")
print(f"可用: {profile.available}")

# 切换策略（必须在 Grounded 状态）
dog.switch_profile("mini")
\end{lstlisting}

\subsection{切换规则}

\begin{itemize}[nosep]
  \item 只能在 \textbf{Grounded} 状态下切换策略
  \item 切换后会重新加载 ONNX 模型
  \item 新策略的 PD 参数和动作缩放会立即生效
\end{itemize}

\section{典型使用流程}

\subsection{标准操作序列}

\begin{lstlisting}[style=python]
import time
from brainstem_sdk import ThunderClient

# 连接
dog = ThunderClient("192.168.66.190")

# 1. 诊断检查
state = dog.get_state()
print(f"初始状态: {state}")          # 应为 "grounded"

voltages = dog.get_voltage()
avg_v = sum(voltages) / len(voltages) if voltages else 0
print(f"电池电压: {avg_v:.1f}V")     # 应 > 42V

motors = dog.get_motor_status()
offline = [m.id for m in motors if not m.online]
if offline:
    print(f"警告: 关节 {offline} 离线")

# 2. 使能电机
dog.enable()

# 3. 站起 (3 秒过渡)
dog.stand_up()
time.sleep(4)
assert dog.get_state() == "standing"

# 4. 行走
dog.walk(vx=0.3)                     # 前进
time.sleep(3)

dog.walk(vx=0.3, vyaw=0.2)          # 前进 + 左转
time.sleep(2)

dog.walk()                            # 停步
time.sleep(1)

# 5. 坐下 (3 秒过渡)
dog.sit_down()
time.sleep(4)
assert dog.get_state() == "grounded"

# 6. 禁用电机
dog.disable()

# 7. 断开连接
dog.close()
\end{lstlisting}

\subsection{实时数据监控}

\begin{lstlisting}[style=python]
import threading
import time
from brainstem_sdk import ThunderClient

dog = ThunderClient("192.168.66.190")

# 在后台线程中监听 IMU
def monitor_imu():
    for imu in dog.listen_imu():
        gx, gy, gz = imu.gyroscope.x, imu.gyroscope.y, imu.gyroscope.z
        print(f"\rIMU: gx={gx:+.2f} gy={gy:+.2f} gz={gz:+.2f}", end="")

imu_thread = threading.Thread(target=monitor_imu, daemon=True)
imu_thread.start()

# 在后台线程中监听状态变化
def monitor_state():
    for state in dog.listen_state():
        print(f"\n[STATE] -> {state}")

state_thread = threading.Thread(target=monitor_state, daemon=True)
state_thread.start()

# 主线程执行操作
dog.enable()
dog.stand_up()
time.sleep(4)
dog.walk(vx=0.3)
time.sleep(5)
dog.sit_down()
time.sleep(4)
dog.disable()
dog.close()
\end{lstlisting}

\subsection{安全包装器}

建议在生产代码中使用 try/finally 确保安全退出：

\begin{lstlisting}[style=python]
import time
from brainstem_sdk import ThunderClient

dog = ThunderClient("192.168.66.190")

try:
    dog.enable()
    dog.stand_up()
    time.sleep(4)

    # ... 您的业务逻辑 ...
    dog.walk(vx=0.3)
    time.sleep(5)

finally:
    # 无论发生什么异常，都确保安全坐下并禁用
    try:
        dog.sit_down()
        time.sleep(5)
    except Exception:
        pass
    try:
        dog.disable()
    except Exception:
        pass
    dog.close()
\end{lstlisting}

\section{关键参数速查}

\begin{longtable}{lll}
\toprule
参数 & 值 & 说明 \\
\midrule
\endhead
控制频率 & 50 Hz (20ms) & 主控制循环和推理频率 \\
gRPC 端口 & 13145 & SDK 连接端口 \\
站起耗时 & 3 秒 (150 ticks) & Grounded $\to$ Standing \\
坐下耗时 & 3 秒 (150 ticks) & Standing $\to$ Grounded \\
仲裁超时 & 3 秒 & 遥控器释放控制权的等待时间 \\
低压阈值 & 42V & 低于此值触发保护坐下 \\
启动超时 & 10 秒 & FSM 等待 Grounded 的超时 \\
关机超时 & 8 秒 & 安全关机等待超时 \\
RPC 超时 & 5 秒 & SDK 默认 RPC 超时 \\
\bottomrule
\end{longtable}



% ============================================================
%  第 5 章：传感器数据
% ============================================================
\chapter{传感器数据}
\label{ch:sensors}

本章介绍 Thunder 四足机器人所有传感器的数据类型、坐标约定、采集方式和编程接口。

\section{传感器总览}

{\small
\begin{longtable}{llllll}
\toprule
传感器 & 硬件 & 数据内容 & 接口方法 & 更新频率 & 传输方式 \\
\midrule
\endhead
IMU & HI91 & 角速度 + 四元数 & \texttt{listen\_imu()} & \textasciitilde 50 Hz & 串口 $\to$ gRPC 流 \\
关节编码器 & Robstride 内置 & 位置/速度/力矩 & \texttt{listen\_joint()} & \textasciitilde 50 Hz & CAN $\to$ gRPC 流 \\
电机状态 & Robstride & 在线/温度/电压/故障 & \texttt{get\_motor\_status()} & 按需查询 & CAN $\to$ gRPC 单次 \\
总线电压 & Robstride & 16 路母线电压 (V) & \texttt{get\_voltage()} & 按需查询 & CAN $\to$ gRPC 单次 \\
状态机 & 软件 & FSM 状态字符串 & \texttt{listen\_state()} & 事件驱动 & gRPC 流 \\
\bottomrule
\end{longtable}
}

\section{IMU 惯性测量单元}

\subsection{硬件与数据}

Thunder 搭载 HI91 惯性测量单元，安装于机身躯干中心，通过串口 (\texttt{/dev/ttyUSB1}) 以约 50 Hz 频率上报数据。每帧包含：

\begin{longtable}{llll}
\toprule
字段 & 类型 & 单位 & 说明 \\
\midrule
\endhead
\texttt{gyroscope.x} & float & rad/s & 绕机体 X 轴角速度 (Roll) \\
\texttt{gyroscope.y} & float & rad/s & 绕机体 Y 轴角速度 (Pitch) \\
\texttt{gyroscope.z} & float & rad/s & 绕机体 Z 轴角速度 (Yaw) \\
\texttt{quaternion.w} & float & 无量纲 & 四元数实部 \\
\texttt{quaternion.x} & float & 无量纲 & 四元数虚部 i \\
\texttt{quaternion.y} & float & 无量纲 & 四元数虚部 j \\
\texttt{quaternion.z} & float & 无量纲 & 四元数虚部 k \\
\texttt{timestamp\_ms} & float & ms & 数据时间戳（毫秒） \\
\bottomrule
\end{longtable}

\subsection{坐标系}

IMU 数据使用\textbf{机体坐标系 (Body Frame)}：

\begin{verbatim}
               +X (前进方向)
                ^
                |
                |
   +Y (左) <---O---> -Y (右)
                |
                |
                v
               -X (后退方向)

   +Z 向上 (垂直地面)
\end{verbatim}

\begin{itemize}[nosep]
  \item X 轴：指向机器人前方
  \item Y 轴：指向机器人左侧
  \item Z 轴：指向上方（右手坐标系）
\end{itemize}

\subsection{四元数约定}

Thunder SDK 使用 \textbf{Hamilton 约定}，四元数顺序为 \texttt{(w, x, y, z)}：

\[
q = w + x\mathbf{i} + y\mathbf{j} + z\mathbf{k}
\]

\begin{itemize}[nosep]
  \item 表示\textbf{世界坐标系到机体坐标系}的旋转
  \item 单位四元数满足 $w^2 + x^2 + y^2 + z^2 = 1$
  \item 机器人水平静止时：$q = (1, 0, 0, 0)$
\end{itemize}

\warning{ROS 等系统常用 \texttt{(x, y, z, w)} 顺序，与本 SDK 不同。如需与 ROS 对接，请手动调整字段顺序。}

\subsection{订阅 IMU 数据}

\begin{lstlisting}[style=python]
from brainstem_sdk import ThunderClient

dog = ThunderClient("192.168.66.190")

# listen_imu() 返回一个迭代器，持续产出 ImuData
for imu in dog.listen_imu():
    gx, gy, gz = imu.gyroscope.x, imu.gyroscope.y, imu.gyroscope.z
    qw, qx, qy, qz = (
        imu.quaternion.w,
        imu.quaternion.x,
        imu.quaternion.y,
        imu.quaternion.z,
    )

    print(f"角速度: ({gx:+.4f}, {gy:+.4f}, {gz:+.4f}) rad/s")
    print(f"四元数: ({qw:.4f}, {qx:.4f}, {qy:.4f}, {qz:.4f})")
    print(f"时间戳: {imu.timestamp_ms:.0f} ms")
\end{lstlisting}

\texttt{listen\_imu()} 是一个阻塞式 gRPC 服务端流，循环会持续运行直到：
\begin{itemize}[nosep]
  \item 手动 \texttt{break} 跳出
  \item 按 \texttt{Ctrl+C} 中断
  \item 网络断开或服务端关闭连接
\end{itemize}

\subsection{数据解读}

\noindent\textbf{角速度解读：}

\begin{longtable}{llll}
\toprule
运动场景 & gyro.x (Roll) & gyro.y (Pitch) & gyro.z (Yaw) \\
\midrule
\endhead
机器人静止 & \textasciitilde 0 & \textasciitilde 0 & \textasciitilde 0 \\
直线行走 & 小幅波动 & 周期性波动 & \textasciitilde 0 \\
原地左转 & \textasciitilde 0 & \textasciitilde 0 & > 0 \\
左侧倾斜 & > 0 & \textasciitilde 0 & \textasciitilde 0 \\
\bottomrule
\end{longtable}

\noindent\textbf{四元数转欧拉角（参考）：}

\begin{lstlisting}[style=python]
import math

def quat_to_euler(q):
    """Hamilton 四元数 (w,x,y,z) 转欧拉角 (roll, pitch, yaw)，单位 rad。"""
    w, x, y, z = q.w, q.x, q.y, q.z
    # Roll (X)
    sinr = 2.0 * (w * x + y * z)
    cosr = 1.0 - 2.0 * (x * x + y * y)
    roll = math.atan2(sinr, cosr)
    # Pitch (Y)
    sinp = 2.0 * (w * y - z * x)
    sinp = max(-1.0, min(1.0, sinp))
    pitch = math.asin(sinp)
    # Yaw (Z)
    siny = 2.0 * (w * z + x * y)
    cosy = 1.0 - 2.0 * (y * y + z * z)
    yaw = math.atan2(siny, cosy)
    return roll, pitch, yaw
\end{lstlisting}

\section{关节编码器}

\subsection{关节布局}

Thunder 共有 16 个自由度，每条腿 3 个旋转关节 + 1 个脚轮：

\begin{verbatim}
  前右 (FR)          前左 (FL)
    0  Hip             4  Hip
    1  Thigh           5  Thigh
    2  Calf            6  Calf
    3  Foot            7  Foot

  后右 (RR)          后左 (RL)
    8  Hip            12  Hip
    9  Thigh          13  Thigh
   10  Calf           14  Calf
   11  Foot           15  Foot
\end{verbatim}

\begin{itemize}[nosep]
  \item \textbf{Hip}：髋关节，控制腿的前后/内外摆动
  \item \textbf{Thigh}：大腿关节，控制大腿抬起/放下
  \item \textbf{Calf}：小腿关节，控制小腿伸展/收缩
  \item \textbf{Foot}：脚轮关节，控制脚部轮毂转动
\end{itemize}

\subsection{数据内容}

每个关节上报以下数据：

\begin{longtable}{llll}
\toprule
字段 & 类型 & 单位 & 说明 \\
\midrule
\endhead
\texttt{id} & int & -- & 关节编号 (0--15) \\
\texttt{position} & float & rad & 当前角位置 \\
\texttt{velocity} & float & rad/s & 当前角速度 \\
\texttt{torque} & float & Nm & 当前力矩 \\
\texttt{status} & int & -- & 状态码 (0 = 正常) \\
\bottomrule
\end{longtable}

\subsection{订阅关节数据}

\begin{lstlisting}[style=python]
from brainstem_sdk import ThunderClient

dog = ThunderClient("192.168.66.190")

print("订阅关节数据 (Ctrl+C 退出)...")
try:
    for joint in dog.listen_joint():
        print(
            f"Joint {joint.id:2d}: "
            f"pos={joint.position:+.4f} rad  "
            f"vel={joint.velocity:+.4f} rad/s  "
            f"torque={joint.torque:+.4f} Nm  "
            f"status={joint.status}"
        )
except KeyboardInterrupt:
    print("\n停止。")
finally:
    dog.close()
\end{lstlisting}

\warning{\texttt{listen\_joint()} 逐关节上报，每次 \texttt{yield} 一个关节的数据。要获取完整的 16 关节快照，需要自行缓存最近一帧的所有关节数据。}

\begin{lstlisting}[style=python]
latest = [None] * 16

for joint in dog.listen_joint():
    latest[joint.id] = joint

    # 检查是否收齐了所有 16 个关节
    if all(j is not None for j in latest):
        # 处理完整快照
        for j in latest:
            print(f"  Joint {j.id}: {j.position:+.3f} rad")
        latest = [None] * 16  # 重置，等待下一帧
\end{lstlisting}

\subsection{关节位置安全限位}

系统默认关节绝对限位为 3.14 rad。任一关节位置的绝对值超过此阈值将立即触发故障保护。该阈值可通过环境变量 \texttt{HAN\_DOG\_JOINT\_LIMIT\_RAD} 调整（设置更小的值可提供提前保护）。

\section{电机状态}

\subsection{查询电机状态}

\texttt{get\_motor\_status()} 返回 16 个 \texttt{MotorInfo} 对象，包含完整的电机健康信息：

\begin{lstlisting}[style=python]
from brainstem_sdk import ThunderClient

dog = ThunderClient("192.168.66.190")

motors = dog.get_motor_status()
for m in motors:
    status = "在线" if m.online else "离线"
    fault = f"故障码: {m.errors}" if m.errors else "无故障"
    print(
        f"Joint {m.id:2d}: {status}  "
        f"温度={m.temperature:.0f}C  "
        f"电压={m.voltage:.1f}V  "
        f"位置={m.position:+.3f}rad  "
        f"速度={m.velocity:+.3f}rad/s  "
        f"力矩={m.torque:+.3f}Nm  "
        f"{fault}"
    )
\end{lstlisting}

\subsection{MotorInfo 字段说明}

\begin{longtable}{llll}
\toprule
字段 & 类型 & 单位 & 说明 \\
\midrule
\endhead
\texttt{id} & int & -- & 关节编号 (0--15) \\
\texttt{online} & bool & -- & 电机是否在线（CAN 通信正常） \\
\texttt{temperature} & float & C & 电机内部温度 \\
\texttt{voltage} & float & V & 电机母线电压 \\
\texttt{position} & float & rad & 当前编码器位置 \\
\texttt{velocity} & float & rad/s & 当前转速 \\
\texttt{torque} & float & Nm & 当前输出力矩 \\
\texttt{errors} & list[int] & -- & 故障码列表，空列表表示无故障 \\
\bottomrule
\end{longtable}

\subsection{故障码解读}

\begin{longtable}{lll}
\toprule
场景 & 表现 & 处理方式 \\
\midrule
\endhead
\texttt{online=False} & 电机未响应 CAN 通信 & 检查 CAN 线缆连接和电源 \\
\texttt{errors} 非空 & 电机报告故障 & 调用 \texttt{clear\_motor\_fault()} 清除 \\
\texttt{temperature > 80} & 电机过热 & 降低负载或暂停运动，等待冷却 \\
\texttt{voltage < 42} & 母线电压过低 & 立即充电，系统将自动触发低压保护 \\
\bottomrule
\end{longtable}

\subsection{清除电机故障}

\begin{lstlisting}[style=python]
# 清除所有电机故障
dog.clear_motor_fault()

# 只清除特定关节（例如前右腿的 3 个关节）
dog.clear_motor_fault(joint_ids=[0, 1, 2])
\end{lstlisting}

\section{电压监控}

\subsection{读取电压}

\texttt{get\_voltage()} 返回 16 个浮点数，对应 16 个电机的母线电压（单位 V）：

\begin{lstlisting}[style=python]
voltages = dog.get_voltage()
print(f"电压范围: {min(voltages):.1f}V ~ {max(voltages):.1f}V")

for i, v in enumerate(voltages):
    warning = " [低压!]" if v < 42.0 else ""
    print(f"  Joint {i:2d}: {v:.1f}V{warning}")
\end{lstlisting}

\subsection{低压保护机制}

\begin{longtable}{ll}
\toprule
阈值 & 行为 \\
\midrule
\endhead
$<$ 42.0 V & 系统自动执行坐下动作，等待 Grounded 状态确认后禁用电机 \\
\bottomrule
\end{longtable}

低压保护是\textbf{不可绕过}的安全机制。当任一电机报告的母线电压低于 42V，控制系统将：

\begin{enumerate}[nosep]
  \item 立即发送 \texttt{SitDown} 指令
  \item 等待机器人进入 Grounded 状态
  \item 禁用所有电机
\end{enumerate}

此过程通过日志输出 \texttt{LOW VOLTAGE: xx.xV < 42.0V -- graceful sitDown}。

\subsection{电压监控最佳实践}

\begin{lstlisting}[style=python]
import time

def monitor_voltage(dog, interval=5.0, warn_threshold=44.0):
    """周期性检查电压，提前预警。"""
    while True:
        voltages = dog.get_voltage()
        min_v = min(voltages)
        if min_v < warn_threshold:
            print(f"[警告] 电压偏低: {min_v:.1f}V (阈值 {warn_threshold}V)")
            if min_v < 42.0:
                print("[严重] 已触发低压保护，请立即充电！")
                break
        else:
            print(f"电压正常: {min_v:.1f}V")
        time.sleep(interval)
\end{lstlisting}

\section{综合示例：传感器数据记录到 CSV}

以下示例同时订阅 IMU 和关节数据，记录到 CSV 文件，便于后续分析：

\begin{lstlisting}[style=python]
"""传感器数据记录器 -- 同时采集 IMU 和关节数据写入 CSV。"""

import csv
import time
import threading
from datetime import datetime
from brainstem_sdk import ThunderClient


def record_sensors(host="192.168.66.190", duration=30.0):
    """记录传感器数据到 CSV 文件。

    Args:
        host: 机器人 IP 地址。
        duration: 记录时长（秒）。
    """
    dog = ThunderClient(host)
    timestamp = datetime.now().strftime("%Y%m%d_%H%M%S")
    stop_event = threading.Event()

    # ---- IMU 记录线程 ----
    def record_imu():
        imu_file = f"imu_{timestamp}.csv"
        with open(imu_file, "w", newline="") as f:
            writer = csv.writer(f)
            writer.writerow([
                "timestamp_ms",
                "gyro_x", "gyro_y", "gyro_z",
                "quat_w", "quat_x", "quat_y", "quat_z",
            ])
            try:
                for imu in dog.listen_imu():
                    if stop_event.is_set():
                        break
                    writer.writerow([
                        f"{imu.timestamp_ms:.1f}",
                        f"{imu.gyroscope.x:.6f}",
                        f"{imu.gyroscope.y:.6f}",
                        f"{imu.gyroscope.z:.6f}",
                        f"{imu.quaternion.w:.6f}",
                        f"{imu.quaternion.x:.6f}",
                        f"{imu.quaternion.y:.6f}",
                        f"{imu.quaternion.z:.6f}",
                    ])
            except Exception as e:
                print(f"IMU 记录异常: {e}")
        print(f"IMU 数据已保存: {imu_file}")

    # ---- 关节记录线程 ----
    def record_joints():
        joint_file = f"joint_{timestamp}.csv"
        with open(joint_file, "w", newline="") as f:
            writer = csv.writer(f)
            writer.writerow([
                "time_s", "joint_id",
                "position_rad", "velocity_rad_s", "torque_nm",
                "status",
            ])
            t0 = time.monotonic()
            try:
                for joint in dog.listen_joint():
                    if stop_event.is_set():
                        break
                    elapsed = time.monotonic() - t0
                    writer.writerow([
                        f"{elapsed:.4f}",
                        joint.id,
                        f"{joint.position:.6f}",
                        f"{joint.velocity:.6f}",
                        f"{joint.torque:.6f}",
                        joint.status,
                    ])
            except Exception as e:
                print(f"关节记录异常: {e}")
        print(f"关节数据已保存: {joint_file}")

    # ---- 启动记录 ----
    imu_thread = threading.Thread(target=record_imu, daemon=True)
    joint_thread = threading.Thread(target=record_joints, daemon=True)

    print(f"开始记录 {duration} 秒...")
    imu_thread.start()
    joint_thread.start()

    try:
        time.sleep(duration)
    except KeyboardInterrupt:
        print("\n手动中断。")

    stop_event.set()
    dog.close()

    imu_thread.join(timeout=3.0)
    joint_thread.join(timeout=3.0)
    print("记录完成。")


if __name__ == "__main__":
    record_sensors(duration=30.0)
\end{lstlisting}

\noindent\textbf{输出文件格式：}

\texttt{imu\_20260403\_143000.csv}:
\begin{lstlisting}[style=plain,numbers=none]
timestamp_ms,gyro_x,gyro_y,gyro_z,quat_w,quat_x,quat_y,quat_z
1234567.0,0.001234,-0.002345,0.000456,0.999800,0.012300,0.005600,0.001200
...
\end{lstlisting}

\texttt{joint\_20260403\_143000.csv}:
\begin{lstlisting}[style=plain,numbers=none]
time_s,joint_id,position_rad,velocity_rad_s,torque_nm,status
0.0000,0,0.123456,0.001234,0.456789,0
0.0001,1,-0.234567,0.002345,-0.567890,0
...
\end{lstlisting}

\section*{附录：传感器数据速查表}

{\small
\begin{longtable}{llllll}
\toprule
数据项 & SDK 方法 & 返回类型 & 单位 & 频率 & 备注 \\
\midrule
\endhead
角速度 & \texttt{listen\_imu()} & \texttt{ImuData.gyroscope} & rad/s & \textasciitilde 50 Hz & 机体坐标系 \\
姿态四元数 & \texttt{listen\_imu()} & \texttt{ImuData.quaternion} & 无量纲 & \textasciitilde 50 Hz & Hamilton (w,x,y,z) \\
IMU 时间戳 & \texttt{listen\_imu()} & \texttt{ImuData.timestamp\_ms} & ms & \textasciitilde 50 Hz & -- \\
关节位置 & \texttt{listen\_joint()} & \texttt{JointData.position} & rad & \textasciitilde 50 Hz & 绝对位置 \\
关节速度 & \texttt{listen\_joint()} & \texttt{JointData.velocity} & rad/s & \textasciitilde 50 Hz & -- \\
关节力矩 & \texttt{listen\_joint()} & \texttt{JointData.torque} & Nm & \textasciitilde 50 Hz & -- \\
电机在线状态 & \texttt{get\_motor\_status()} & \texttt{MotorInfo.online} & bool & 按需 & CAN 通信 \\
电机温度 & \texttt{get\_motor\_status()} & \texttt{MotorInfo.temperature} & C & 按需 & 内部温度 \\
电机电压 & \texttt{get\_motor\_status()} & \texttt{MotorInfo.voltage} & V & 按需 & 母线电压 \\
电机故障码 & \texttt{get\_motor\_status()} & \texttt{MotorInfo.errors} & list[int] & 按需 & 空=无故障 \\
母线电压(16路) & \texttt{get\_voltage()} & list[float] & V & 按需 & 低压阈值 42V \\
FSM 状态变化 & \texttt{listen\_state()} & str & -- & 事件驱动 & 状态切换时触发 \\
\bottomrule
\end{longtable}
}



% ============================================================
%  第 6 章：API 参考
% ============================================================
\chapter{API 参考}
\label{ch:api}

本章提供 Thunder Python SDK (\texttt{brainstem\_sdk}) 全部公开接口的完整参考。

\section{ThunderClient 构造与连接}

\subsection{构造函数}

\begin{lstlisting}[style=python]
ThunderClient(host="192.168.66.190", port=13145, timeout=5.0)
\end{lstlisting}

\begin{longtable}{llll}
\toprule
参数 & 类型 & 默认值 & 说明 \\
\midrule
\endhead
\texttt{host} & str & \texttt{"192.168.66.190"} & 机器人 IP 地址 \\
\texttt{port} & int & \texttt{13145} & gRPC 服务端口 \\
\texttt{timeout} & float & \texttt{5.0} & 默认 RPC 超时时间（秒） \\
\bottomrule
\end{longtable}

\noindent\textbf{说明}：构造函数创建一个 gRPC 不安全通道（insecure channel），连接到 \texttt{host:port}。通道创建即连接，无需额外调用。

\begin{lstlisting}[style=python]
from brainstem_sdk import ThunderClient

# 默认连接
dog = ThunderClient("192.168.66.190")

# 自定义参数
dog = ThunderClient("10.0.0.100", port=13145, timeout=10.0)
\end{lstlisting}

\subsection{close()}

\begin{lstlisting}[style=python,numbers=none]
dog.close() -> None
\end{lstlisting}

关闭 gRPC 通道，释放网络资源。关闭后不可再调用任何方法。

\subsection{上下文管理器}

\texttt{ThunderClient} 支持 \texttt{with} 语句，退出时自动调用 \texttt{close()}：

\begin{lstlisting}[style=python]
with ThunderClient("192.168.66.190") as dog:
    dog.enable()
    dog.stand_up()
    dog.walk(vx=0.3)
    dog.sit_down()
    dog.disable()
# 退出 with 块后通道自动关闭
\end{lstlisting}

\section{电机控制}

\subsection{enable()}

\begin{lstlisting}[style=python,numbers=none]
dog.enable() -> None
\end{lstlisting}

使能全部 16 个电机。调用后电机通电并进入闭环控制。

\begin{longtable}{ll}
\toprule
项目 & 说明 \\
\midrule
\endhead
前置条件 & gRPC 连接正常，控制权未被遥控器占用 \\
异常 & \texttt{grpc.RpcError} -- 超时或连接失败 \\
\bottomrule
\end{longtable}

\subsection{disable()}

\begin{lstlisting}[style=python,numbers=none]
dog.disable() -> None
\end{lstlisting}

禁用全部电机（安全停止）。电机断电，关节自由落体。

\begin{longtable}{ll}
\toprule
项目 & 说明 \\
\midrule
\endhead
前置条件 & 无（任何状态下均可调用） \\
建议 & 在 Grounded 状态下调用，避免机器人从站立姿态跌落 \\
异常 & \texttt{grpc.RpcError} -- 超时或连接失败 \\
\bottomrule
\end{longtable}

\begin{lstlisting}[style=python]
# 安全关机序列
dog.walk(vx=0.0)   # 停步
dog.sit_down()      # 坐下
# 等待机器人坐稳...
dog.disable()       # 禁用电机
\end{lstlisting}

\section{运动指令}

\subsection{walk()}

\begin{lstlisting}[style=python,numbers=none]
dog.walk(vx=0.0, vy=0.0, vyaw=0.0) -> None
\end{lstlisting}

发送行走速度指令。

\begin{longtable}{lllll}
\toprule
参数 & 类型 & 默认值 & 范围 & 说明 \\
\midrule
\endhead
\texttt{vx} & float & \texttt{0.0} & {[}-1.0, 1.0{]} & 前后速度。正值=前进，负值=后退 \\
\texttt{vy} & float & \texttt{0.0} & {[}-1.0, 1.0{]} & 左右速度。正值=向左，负值=向右 \\
\texttt{vyaw} & float & \texttt{0.0} & {[}-1.0, 1.0{]} & 旋转速度。正值=逆时针，负值=顺时针 \\
\bottomrule
\end{longtable}

\begin{longtable}{ll}
\toprule
项目 & 说明 \\
\midrule
\endhead
前置条件 & 机器人处于 Standing 或 Walking 状态 \\
停步 & 调用 \texttt{dog.walk()} 或 \texttt{dog.walk(0, 0, 0)} \\
异常 & \texttt{grpc.RpcError} -- 超时、连接失败或 Arbiter 拒绝 \\
\bottomrule
\end{longtable}

\noindent\textbf{速度参数为归一化值}，不是物理速度。\texttt{vx=1.0} 表示当前策略允许的最大前进速度，实际物理速度取决于所加载的运动策略。

\begin{lstlisting}[style=python]
# 前进
dog.walk(vx=0.5)

# 向左平移
dog.walk(vy=0.3)

# 前进 + 左转
dog.walk(vx=0.3, vyaw=0.2)

# 原地停步
dog.walk()
\end{lstlisting}

\subsection{stand\_up()}

\begin{lstlisting}[style=python,numbers=none]
dog.stand_up() -> None
\end{lstlisting}

从坐姿过渡到站立。

\begin{longtable}{ll}
\toprule
项目 & 说明 \\
\midrule
\endhead
前置条件 & 机器人处于 Grounded 状态，电机已使能 \\
状态变化 & Grounded $\to$ Transitioning $\to$ Standing \\
异常 & \texttt{grpc.RpcError} -- 超时或连接失败 \\
\bottomrule
\end{longtable}

\subsection{sit\_down()}

\begin{lstlisting}[style=python,numbers=none]
dog.sit_down() -> None
\end{lstlisting}

从站立过渡到坐姿。如果正在行走，会先停步再坐下。

\begin{longtable}{ll}
\toprule
项目 & 说明 \\
\midrule
\endhead
前置条件 & 机器人处于 Standing 或 Walking 状态 \\
状态变化 & Standing/Walking $\to$ Transitioning $\to$ Grounded \\
异常 & \texttt{grpc.RpcError} -- 超时或连接失败 \\
\bottomrule
\end{longtable}

\section{状态查询}

\subsection{get\_state()}

\begin{lstlisting}[style=python,numbers=none]
dog.get_state() -> str
\end{lstlisting}

查询当前机器人状态。

\begin{longtable}{ll}
\toprule
返回值 & 说明 \\
\midrule
\endhead
\texttt{"zero"} & 初始化 / 未使能 \\
\texttt{"grounded"} & 坐姿（安全状态） \\
\texttt{"standing"} & 站立，可接收行走指令 \\
\texttt{"walking"} & 行走中 \\
\texttt{"transitioning"} & 状态过渡中（站起 / 坐下） \\
\bottomrule
\end{longtable}

\begin{lstlisting}[style=python]
state = dog.get_state()
print(f"当前状态: {state}")

if state == "standing":
    dog.walk(vx=0.3)
\end{lstlisting}

\subsection{get\_voltage()}

\begin{lstlisting}[style=python,numbers=none]
dog.get_voltage() -> list[float]
\end{lstlisting}

读取 16 个电机的母线电压。

\begin{longtable}{ll}
\toprule
项目 & 说明 \\
\midrule
\endhead
返回值 & 长度为 16 的浮点数列表，单位 V \\
索引 & 对应关节编号 0--15 \\
低压阈值 & 42.0V，低于此值系统自动触发保护 \\
\bottomrule
\end{longtable}

\subsection{get\_motor\_status()}

\begin{lstlisting}[style=python,numbers=none]
dog.get_motor_status() -> list[MotorInfo]
\end{lstlisting}

获取全部 16 个电机的健康状态。

\begin{lstlisting}[style=python]
motors = dog.get_motor_status()
for m in motors:
    if not m.online:
        print(f"Joint {m.id} 离线!")
    if m.errors:
        print(f"Joint {m.id} 故障: {m.errors}")
    if m.temperature > 70:
        print(f"Joint {m.id} 温度偏高: {m.temperature}C")
\end{lstlisting}

\section{策略管理}

\subsection{get\_profile()}

\begin{lstlisting}[style=python,numbers=none]
dog.get_profile() -> ProfileInfo
\end{lstlisting}

获取当前运动策略信息及可用策略列表。

\begin{lstlisting}[style=python]
profile = dog.get_profile()
print(f"当前策略: {profile.current}")
print(f"可用策略: {profile.available}")
for name, desc in zip(profile.available, profile.descriptions):
    print(f"  {name}: {desc}")
\end{lstlisting}

\subsection{switch\_profile()}

\begin{lstlisting}[style=python,numbers=none]
dog.switch_profile(name: str) -> ProfileInfo
\end{lstlisting}

切换运动策略。

\begin{longtable}{lll}
\toprule
参数 & 类型 & 说明 \\
\midrule
\endhead
\texttt{name} & str & 目标策略名称，必须在 \texttt{get\_profile().available} 列表中 \\
\bottomrule
\end{longtable}

\begin{longtable}{ll}
\toprule
项目 & 说明 \\
\midrule
\endhead
前置条件 & 机器人处于 \textbf{Grounded} 状态 \\
返回值 & 切换后的 \texttt{ProfileInfo} \\
异常 & \texttt{grpc.RpcError} -- 策略不存在、状态不对或超时 \\
\bottomrule
\end{longtable}

\begin{lstlisting}[style=python]
# 查看可用策略
profile = dog.get_profile()
print(f"可用: {profile.available}")

# 切换策略（必须先坐下）
dog.sit_down()
# 等待 Grounded ...
new_profile = dog.switch_profile("thunder2v1")
print(f"已切换到: {new_profile.current}")
\end{lstlisting}

\section{诊断与标定}

\subsection{clear\_motor\_fault()}

\begin{lstlisting}[style=python,numbers=none]
dog.clear_motor_fault(joint_ids=None) -> None
\end{lstlisting}

清除电机故障码。

\begin{longtable}{llll}
\toprule
参数 & 类型 & 默认值 & 说明 \\
\midrule
\endhead
\texttt{joint\_ids} & list[int] | None & \texttt{None} & 要清除的关节编号列表 (0--15)。\texttt{None} 表示清除全部 \\
\bottomrule
\end{longtable}

\begin{lstlisting}[style=python]
# 清除所有电机故障
dog.clear_motor_fault()

# 只清除前右腿
dog.clear_motor_fault(joint_ids=[0, 1, 2, 3])

# 只清除单个关节
dog.clear_motor_fault(joint_ids=[5])
\end{lstlisting}

\subsection{set\_zero()}

\begin{lstlisting}[style=python,numbers=none]
dog.set_zero() -> None
\end{lstlisting}

将当前关节位置标定为零位。

\begin{longtable}{ll}
\toprule
项目 & 说明 \\
\midrule
\endhead
前置条件 & 机器人处于 \textbf{Grounded} 状态 \\
用途 & 出厂标定或维修后重新校准零位 \\
注意 & 错误的零位会导致运动异常，非必要不调用 \\
\bottomrule
\end{longtable}

\begin{lstlisting}[style=python]
# 确保在 Grounded 状态下执行
state = dog.get_state()
if state == "grounded":
    dog.set_zero()
    print("零位标定完成。")
else:
    print(f"当前状态 {state}，请先坐下。")
\end{lstlisting}

\section{实时数据流}

所有 \texttt{listen\_*} 方法返回 Python 迭代器（\texttt{Iterator}），底层为 gRPC 服务端流。迭代器在以下情况终止：

\begin{itemize}[nosep]
  \item 手动 \texttt{break}
  \item 键盘中断 (\texttt{Ctrl+C})
  \item 调用 \texttt{dog.close()} 关闭通道
  \item 网络断开
\end{itemize}

\subsection{listen\_imu()}

\begin{lstlisting}[style=python,numbers=none]
dog.listen_imu() -> Iterator[ImuData]
\end{lstlisting}

订阅 IMU 实时数据流，约 50Hz。

\begin{longtable}{ll}
\toprule
项目 & 说明 \\
\midrule
\endhead
产出 & \texttt{ImuData}（含 \texttt{gyroscope}, \texttt{quaternion}, \texttt{timestamp\_ms}） \\
频率 & \textasciitilde 50 Hz \\
阻塞 & 是（在 \texttt{for} 循环中阻塞等待下一帧） \\
\bottomrule
\end{longtable}

\subsection{listen\_joint()}

\begin{lstlisting}[style=python,numbers=none]
dog.listen_joint() -> Iterator[JointData]
\end{lstlisting}

订阅关节实时数据流。

\begin{longtable}{ll}
\toprule
项目 & 说明 \\
\midrule
\endhead
产出 & \texttt{JointData}（含 \texttt{id}, \texttt{position}, \texttt{velocity}, \texttt{torque}, \texttt{status}） \\
频率 & \textasciitilde 50 Hz（逐关节上报） \\
阻塞 & 是 \\
\bottomrule
\end{longtable}

\subsection{listen\_state()}

\begin{lstlisting}[style=python,numbers=none]
dog.listen_state() -> Iterator[str]
\end{lstlisting}

订阅状态机变化事件。仅在状态切换时产出，不定频。

\begin{longtable}{ll}
\toprule
项目 & 说明 \\
\midrule
\endhead
产出 & 状态字符串：\texttt{"zero"}, \texttt{"grounded"}, \texttt{"standing"}, \texttt{"walking"}, \texttt{"transitioning"} \\
频率 & 事件驱动（仅状态变化时触发） \\
阻塞 & 是 \\
\bottomrule
\end{longtable}

\section{数据类型参考}

\subsection{Vec3}

三维向量。

\begin{lstlisting}[style=python]
@dataclass
class Vec3:
    x: float = 0.0
    y: float = 0.0
    z: float = 0.0
\end{lstlisting}

\subsection{Quat}

四元数（Hamilton 约定：w, x, y, z）。

\begin{lstlisting}[style=python]
@dataclass
class Quat:
    w: float = 1.0    # 实部
    x: float = 0.0    # 虚部 i
    y: float = 0.0    # 虚部 j
    z: float = 0.0    # 虚部 k
\end{lstlisting}

\subsection{ImuData}

IMU 惯性测量数据帧。

\begin{lstlisting}[style=python]
@dataclass
class ImuData:
    gyroscope: Vec3       # 角速度 (rad/s)
    quaternion: Quat      # 姿态四元数 (Hamilton w,x,y,z)
    timestamp_ms: float   # 时间戳 (ms)
\end{lstlisting}

\subsection{JointData}

单个关节的实时数据。

\begin{lstlisting}[style=python]
@dataclass
class JointData:
    id: int            # 关节编号 (0-15)
    position: float    # 角位置 (rad)
    velocity: float    # 角速度 (rad/s)
    torque: float      # 力矩 (Nm)
    status: int        # 状态码 (0=正常)
\end{lstlisting}

\subsection{MotorInfo}

电机完整健康状态。

\begin{lstlisting}[style=python]
@dataclass
class MotorInfo:
    id: int              # 关节编号 (0-15)
    online: bool         # 是否在线
    temperature: float   # 温度 (C)
    voltage: float       # 母线电压 (V)
    position: float      # 编码器位置 (rad)
    velocity: float      # 转速 (rad/s)
    torque: float        # 输出力矩 (Nm)
    errors: list[int]    # 故障码列表 (空=无故障)
\end{lstlisting}

\subsection{ProfileInfo}

运动策略信息。

\begin{lstlisting}[style=python]
@dataclass
class ProfileInfo:
    current: str           # 当前激活的策略名称
    available: list[str]   # 所有可用策略名称
    descriptions: list[str] # 各策略说明文字
\end{lstlisting}

\subsection{RobotState}

状态机常量（字符串）。

\begin{lstlisting}[style=python]
class RobotState:
    ZERO = "zero"
    GROUNDED = "grounded"
    STANDING = "standing"
    WALKING = "walking"
    TRANSITIONING = "transitioning"
\end{lstlisting}

使用方式：

\begin{lstlisting}[style=python]
from brainstem_sdk import RobotState

state = dog.get_state()
if state == RobotState.STANDING:
    dog.walk(vx=0.3)
\end{lstlisting}

\section{错误处理}

所有 RPC 调用可能抛出 \texttt{grpc.RpcError}。以下是常见错误场景及处理方式：

\subsection{连接超时}

\begin{lstlisting}[style=python]
import grpc

try:
    dog = ThunderClient("192.168.66.190", timeout=3.0)
    dog.get_state()
except grpc.RpcError as e:
    if e.code() == grpc.StatusCode.UNAVAILABLE:
        print("无法连接到机器人。请检查：")
        print("  1. 机器人是否已开机")
        print("  2. IP 地址是否正确")
        print("  3. brainstem 服务是否正在运行")
    elif e.code() == grpc.StatusCode.DEADLINE_EXCEEDED:
        print("请求超时。网络可能不稳定。")
    else:
        print(f"gRPC 错误: {e.code()} - {e.details()}")
\end{lstlisting}

\subsection{常见错误码}

\begin{longtable}{lll}
\toprule
gRPC 状态码 & 含义 & 常见原因 \\
\midrule
\endhead
\texttt{UNAVAILABLE} & 服务不可达 & 机器人未开机、IP 错误、服务未启动 \\
\texttt{DEADLINE\_EXCEEDED} & 请求超时 & 网络延迟、服务端无响应 \\
\texttt{FAILED\_PRECONDITION} & 前置条件不满足 & 状态不对（如在 Zero 状态下调用 walk） \\
\texttt{PERMISSION\_DENIED} & 权限被拒 & Arbiter 控制权被遥控器占用 \\
\texttt{INTERNAL} & 服务端内部错误 & 硬件异常或程序错误 \\
\bottomrule
\end{longtable}

\subsection{Arbiter 控制权冲突}

Thunder 的控制仲裁器（ControlArbiter）确保遥控器始终具有最高优先级。当遥控器在线时，gRPC 指令可能被拒绝：

\begin{itemize}[nosep]
  \item 遥控器在线且正在操作：gRPC 运动指令将被忽略
  \item 遥控器在线但空闲超过 3 秒：gRPC 自动获得控制权
  \item 遥控器离线：gRPC 始终拥有控制权
\end{itemize}

\begin{lstlisting}[style=python]
import grpc
import time

def safe_walk(dog, vx=0.0, vy=0.0, vyaw=0.0, retries=3):
    """带重试的行走指令。"""
    for attempt in range(retries):
        try:
            dog.walk(vx=vx, vy=vy, vyaw=vyaw)
            return True
        except grpc.RpcError as e:
            if e.code() == grpc.StatusCode.PERMISSION_DENIED:
                print(f"控制权被遥控器占用，等待 ({attempt+1}/{retries})...")
                time.sleep(1.0)
            else:
                raise
    print("无法获得控制权。请确认遥控器已松开操控杆。")
    return False
\end{lstlisting}

\subsection{流式接口异常处理}

\begin{lstlisting}[style=python]
try:
    for imu in dog.listen_imu():
        # 处理数据...
        pass
except grpc.RpcError as e:
    if e.code() == grpc.StatusCode.UNAVAILABLE:
        print("连接断开，尝试重连...")
    else:
        print(f"流异常: {e.code()}")
except KeyboardInterrupt:
    print("用户中断。")
finally:
    dog.close()
\end{lstlisting}

\section*{附录：方法速查表}

{\small
\begin{longtable}{llll}
\toprule
分类 & 方法 & 返回值 & 说明 \\
\midrule
\endhead
\textbf{连接} & \texttt{ThunderClient(host, port, timeout)} & -- & 构造并连接 \\
 & \texttt{close()} & None & 关闭连接 \\
\textbf{电机} & \texttt{enable()} & None & 使能全部电机 \\
 & \texttt{disable()} & None & 禁用全部电机 \\
\textbf{运动} & \texttt{walk(vx, vy, vyaw)} & None & 行走指令 \\
 & \texttt{stand\_up()} & None & 坐姿 $\to$ 站立 \\
 & \texttt{sit\_down()} & None & 站立 $\to$ 坐姿 \\
\textbf{状态} & \texttt{get\_state()} & str & 查询 FSM 状态 \\
 & \texttt{get\_voltage()} & list[float] & 16 路母线电压 \\
 & \texttt{get\_motor\_status()} & list[MotorInfo] & 16 个电机健康状态 \\
\textbf{策略} & \texttt{get\_profile()} & ProfileInfo & 当前策略信息 \\
 & \texttt{switch\_profile(name)} & ProfileInfo & 切换运动策略 \\
\textbf{诊断} & \texttt{clear\_motor\_fault(joint\_ids)} & None & 清除电机故障 \\
 & \texttt{set\_zero()} & None & 标定零位 \\
\textbf{数据流} & \texttt{listen\_imu()} & Iterator[ImuData] & IMU 实时流 \\
 & \texttt{listen\_joint()} & Iterator[JointData] & 关节实时流 \\
 & \texttt{listen\_state()} & Iterator[str] & 状态变化流 \\
\bottomrule
\end{longtable}
}



% ============================================================
%  第 7 章：安全指南
% ============================================================
\chapter{安全指南}
\label{ch:safety}

Thunder 四足机器人是一台具有高扭矩电机的自主运动设备。操作不当可能造成设备损坏或人员伤害。请在使用前仔细阅读本章全部内容。

\section{操作前检查清单}

\textbf{每次启动机器人前，必须逐项确认以下事项：}

{\small
\begin{longtable}{clll}
\toprule
序号 & 检查项 & 确认方法 & 通过标准 \\
\midrule
\endhead
1 & 电池电量充足 & \texttt{dog.get\_voltage()} & 所有电机电压 >= 44V \\
2 & 操作区域净空 & 目视检查 & 周围 1.5 米内无人员/易碎物品 \\
3 & 机器人放置在地面 & 目视检查 & 四脚均接触平坦地面 \\
4 & 急停装置就绪 & 确认遥控器在手 & 遥控器已开机，电量充足 \\
5 & CAN 总线连接牢固 & 目视检查 & 所有 CAN 线缆插头无松动 \\
6 & 控制服务已启动 & SSH 检查 & \texttt{systemctl status brainstem} 显示 active \\
7 & 电机状态正常 & \texttt{get\_motor\_status()} & 16 个电机全部在线，无故障 \\
8 & 关节位置合理 & \texttt{get\_motor\_status()} & 所有关节位置绝对值 $<$ 1.0 rad \\
9 & 遥控器优先级确认 & 操控遥控器摇杆 & 遥控器能正常控制机器人 \\
10 & 运行环境温度适宜 & 环境检查 & 0--40\textdegree C 室温，非湿滑地面 \\
\bottomrule
\end{longtable}
}

\noindent\textbf{完整的启动前检查代码：}

\begin{lstlisting}[style=python]
from brainstem_sdk import ThunderClient

dog = ThunderClient("192.168.66.190")

# 1. 检查电压
voltages = dog.get_voltage()
min_v = min(voltages)
print(f"电压: {min_v:.1f}V", "-- 通过" if min_v >= 44.0 else "-- 不通过! 请充电")

# 2. 检查电机状态
motors = dog.get_motor_status()
offline = [m for m in motors if not m.online]
faulted = [m for m in motors if m.errors]

if offline:
    print(f"离线电机: {[m.id for m in offline]} -- 不通过!")
else:
    print("电机在线: 全部 16/16 -- 通过")

if faulted:
    print(f"故障电机: {[m.id for m in faulted]} -- 尝试清除...")
    dog.clear_motor_fault()
else:
    print("电机故障: 无 -- 通过")

# 3. 检查关节位置
bad_joints = [m for m in motors if abs(m.position) > 1.0]
if bad_joints:
    print(f"关节位置异常: {[(m.id, f'{m.position:.2f}rad') for m in bad_joints]}")
else:
    print("关节位置: 正常 -- 通过")

# 4. 检查温度
hot = [m for m in motors if m.temperature > 60]
if hot:
    print(f"温度偏高: {[(m.id, f'{m.temperature:.0f}C') for m in hot]}")
else:
    print("电机温度: 正常 -- 通过")

# 5. 检查状态
state = dog.get_state()
print(f"当前状态: {state}", "-- 通过" if state == "grounded" else "-- 注意: 非 Grounded")

dog.close()
\end{lstlisting}

\section{低压保护}

\subsection{保护机制}

当任一电机报告的母线电压低于 \textbf{42.0V} 时，系统自动触发低压保护：

\begin{verbatim}
检测到低压 (< 42V)
       |
       v
  发送 SitDown 指令
       |
       v
  等待进入 Grounded 状态
       |
       v
  禁用全部电机 (Disable)
       |
       v
  日志输出: "LOW VOLTAGE: xx.xV < 42.0V -- graceful sitDown"
\end{verbatim}

此保护是\textbf{优雅坐下}流程，不是直接断电。系统会等待机器人安全着地后才禁用电机，避免从站立姿态跌落。

\subsection{电压等级参考}

\begin{longtable}{lll}
\toprule
电压范围 & 状态 & 建议 \\
\midrule
\endhead
$\geq$ 48V & 满电 & 正常使用 \\
44V -- 48V & 正常 & 正常使用，关注电量 \\
42V -- 44V & 低电量 & 尽快结束任务，准备充电 \\
$<$ 42V & 危险 & 自动触发保护，必须立即充电 \\
\bottomrule
\end{longtable}

\subsection{注意事项}

\begin{itemize}[nosep]
  \item 低压保护\textbf{不可通过 SDK 绕过}
  \item 大电流运动（快速行走、爬坡）会导致瞬时压降，可能在电池未耗尽时触发保护
  \item 建议在电压低于 44V 时即结束运动任务
  \item 触发低压保护后，必须充电至 44V 以上才可重新启动
\end{itemize}

\section{关节限位保护}

\subsection{保护机制}

每个关节的位置都受到绝对限位保护。默认限位为 \textbf{3.14 rad}（即 $\pi$，约 180 度）。任一关节位置绝对值超过此阈值将立即触发故障（Fault）。

\begin{longtable}{llll}
\toprule
配置 & 环境变量 & 默认值 & 说明 \\
\midrule
\endhead
关节限位 & \texttt{HAN\_DOG\_JOINT\_LIMIT\_RAD} & 3.14 & 绝对值限位，单位 rad \\
\bottomrule
\end{longtable}

\subsection{触发后果}

关节超限后：
\begin{enumerate}[nosep]
  \item 对应电机立即进入故障状态
  \item 机器人可能失去该腿的控制，导致摔倒
  \item 必须手动清除故障后才能恢复
\end{enumerate}

\subsection{预防措施}

\begin{itemize}[nosep]
  \item 确保零位标定正确（\texttt{set\_zero()} 在标准坐姿下执行）
  \item 不要在关节卡住或受阻的情况下强行运动
  \item 如果需要更严格的保护，可设置更小的限位值：
\end{itemize}

\begin{lstlisting}[style=bash]
export HAN_DOG_JOINT_LIMIT_RAD=2.5  # 约 143 度
\end{lstlisting}

\section{电机故障处理}

\subsection{检测故障}

\begin{lstlisting}[style=python]
motors = dog.get_motor_status()
for m in motors:
    if m.errors:
        print(f"Joint {m.id} 故障码: {m.errors}")
    if not m.online:
        print(f"Joint {m.id} 离线 (CAN 通信中断)")
\end{lstlisting}

\subsection{清除故障}

\begin{lstlisting}[style=python]
# 方法 1: 清除所有电机故障
dog.clear_motor_fault()

# 方法 2: 只清除特定关节
dog.clear_motor_fault(joint_ids=[0, 1, 2])
\end{lstlisting}

\subsection{故障处理流程}

\begin{verbatim}
发现电机故障
       |
       v
  机器人是否在运动中?
       |
    是/        \否
    v            v
  立即坐下     直接清除故障
  dog.sit_down()    dog.clear_motor_fault()
       |            |
       v            v
  等待 Grounded   重新检查
       |            |
       v            v
  清除故障       故障清除?
  dog.clear_motor_fault()   |
       |          是/    \否
       v          v        v
  重新检查     恢复运动   检查硬件
                          (CAN线/电源/电机)
\end{verbatim}

\subsection{反复出现的故障}

如果 \texttt{clear\_motor\_fault()} 后故障立即重新出现，说明存在持续性硬件问题：

\begin{longtable}{lll}
\toprule
表现 & 可能原因 & 处理 \\
\midrule
\endhead
单个电机反复故障 & CAN 线缆松动 & 重新插拔对应 CAN 接头 \\
多个电机同时故障 & 母线问题 & 检查电池和配电板 \\
故障码伴随高温 & 电机过热 & 停止运动，等待冷却至 50\textdegree C 以下 \\
关节位置跳变后故障 & 编码器异常 & 重新标零 \texttt{set\_zero()} \\
\bottomrule
\end{longtable}

\section{急停程序}

按紧急程度从低到高，有三种急停方式：

\subsection{软件急停（SDK）}

适用于程序可控的紧急情况：

\begin{lstlisting}[style=python]
# 优雅急停：停步 -> 坐下 -> 禁用
dog.walk(vx=0.0, vy=0.0, vyaw=0.0)
dog.sit_down()
# 等待进入 Grounded ...
dog.disable()
\end{lstlisting}

\subsection{遥控器急停}

适用于现场操作人员发现危险时：

\begin{enumerate}[nosep]
  \item \textbf{立即操控遥控器摇杆} -- 遥控器自动获得最高控制权，接管机器人
  \item 通过遥控器操作坐下
  \item 遥控器上的急停按钮可直接断开电机使能
\end{enumerate}

\warning{遥控器始终具有最高优先级，可在任何时刻覆盖 SDK 指令。}

\subsection{硬件断电}

适用于最紧急的情况（机器人失控、硬件冒烟等）：

\begin{enumerate}[nosep]
  \item \textbf{直接关闭电池开关或拔出电池}
  \item 机器人将立即断电跌落
  \item 注意：断电跌落可能造成机械损伤，仅在绝对必要时使用
\end{enumerate}

\subsection{急停优先级}

\begin{verbatim}
优先级从高到低:

  [1] 硬件断电       -- 物理级，不可覆盖
  [2] 遥控器控制     -- Arbiter 保证最高软件优先级
  [3] SDK disable()  -- gRPC 调用
  [4] 低压自动保护   -- 系统自动触发
\end{verbatim}

\section{Arbiter 控制仲裁安全}

\subsection{仲裁规则}

Thunder 的 ControlArbiter（控制仲裁器）实施以下安全规则：

\begin{longtable}{ll}
\toprule
规则 & 说明 \\
\midrule
\endhead
遥控器优先 & YUNZHUO 遥控器的优先级永远高于 gRPC（SDK） \\
超时释放 & 遥控器停止操作 3 秒后，gRPC 自动获得控制权 \\
无抢占 & gRPC 不能强行抢占遥控器的控制权 \\
超时可配 & 超时时间通过 \texttt{HAN\_DOG\_ARBITER\_TIMEOUT} 配置，默认 3 秒 \\
\bottomrule
\end{longtable}

\subsection{对 SDK 开发者的影响}

\begin{verbatim}
SDK 发送 walk() 指令
       |
       v
  Arbiter 检查控制权
       |
   gRPC 有控制权?
     |
   是/      \否 (遥控器占用)
   v          v
 执行指令    指令被忽略或拒绝
\end{verbatim}

\begin{itemize}[nosep]
  \item 当遥控器在线并活跃时，SDK 的运动指令将被忽略
  \item 状态查询（\texttt{get\_state()}, \texttt{get\_voltage()} 等）不受 Arbiter 影响，始终可用
  \item 遥控器释放后约 3 秒，SDK 自动恢复控制权，无需额外操作
\end{itemize}

\subsection{安全启示}

\begin{itemize}[nosep]
  \item \textbf{现场测试时务必携带遥控器}，作为最后的安全手段
  \item 不要依赖 SDK 作为唯一的急停方式
  \item 如果 SDK 指令无响应，先检查遥控器是否在线
\end{itemize}

\section{禁止操作}

以下操作可能导致设备损坏或人员伤害，\textbf{严格禁止}：

{\small
\begin{longtable}{cll}
\toprule
序号 & 禁止操作 & 原因 \\
\midrule
\endhead
1 & 在桌面、高台上运行机器人 & 摔落损坏，可能伤人 \\
2 & 电机使能状态下翻转机器人 & 电机强力对抗将损坏齿轮箱 \\
3 & 跳过 Grounded 状态直接启用电机运动 & 状态机保护被绕过 \\
4 & 在电机使能状态下插拔 CAN 线缆 & 可能导致 CAN 总线短路 \\
5 & 长时间运行而不监控温度 & 电机过热会永久损坏绕组 \\
6 & 在关节卡死时反复发送运动指令 & 电流持续上升，烧毁电机 \\
7 & 在充电状态下运行电机 & 部分充电器不支持边充边放 \\
8 & 修改关节限位为极大值 & 失去关节保护 \\
9 & 在湿滑/不平地面运行而无人看护 & 跌倒后持续运动可能损坏关节 \\
10 & 多个 SDK 客户端同时发送运动指令 & 指令冲突导致运动抖动 \\
\bottomrule
\end{longtable}
}

\section{标准关机序列}

每次使用结束后，必须执行以下关机序列：

\begin{lstlisting}[style=python]
import time

# 步骤 1: 停止行走
dog.walk(vx=0.0, vy=0.0, vyaw=0.0)

# 步骤 2: 坐下
dog.sit_down()

# 步骤 3: 等待进入 Grounded 状态
for state in dog.listen_state():
    if state == "grounded":
        break

# 步骤 4: 禁用电机
dog.disable()

# 步骤 5: 关闭连接
dog.close()

print("安全关机完成。")
\end{lstlisting}

\noindent\textbf{简化版（带超时保护）：}

\begin{lstlisting}[style=python]
import time

def safe_shutdown(dog, timeout=10.0):
    """安全关机，带超时保护。"""
    try:
        dog.walk()
        dog.sit_down()

        # 轮询等待 Grounded，而非使用流式接口
        deadline = time.monotonic() + timeout
        while time.monotonic() < deadline:
            if dog.get_state() == "grounded":
                break
            time.sleep(0.2)

        dog.disable()
    except Exception as e:
        print(f"关机异常: {e}")
        # 最后手段：直接禁用
        try:
            dog.disable()
        except Exception:
            pass
    finally:
        dog.close()
\end{lstlisting}

\noindent\textbf{关键原则}：先坐稳再断电。直接调用 \texttt{disable()} 而不先 \texttt{sit\_down()} 会导致机器人从站立姿态自由跌落。

\section{故障诊断树}

遇到问题时，按以下诊断树逐步排查：

\begin{verbatim}
问题: 机器人异常
       |
       v
  机器人是否有响应?
       |
    否/        \是
    v            v
  检查电源      机器人在运动吗?
  检查网络        |
  重启服务      是/        \否
       |       v            v
       |    运动是否正常?   能查询状态吗?
       |      |              dog.get_state()
       |    否/  \是           |
       |    v      v        是/   \否
       |  运动异常  正常     查电机状态  网络问题
       |         get_motor_status()
       v                     |
  [电源/网络故障]         有故障码?
  1. 电池是否有电           |
  2. 电池开关是否打开    是/    \否
  3. CAN 线是否连好      v        v
  4. SSH 能否连通      清除故障  检查状态机
  5. brainstem 服务    clear_motor_fault()
     是否运行
\end{verbatim}

\noindent\textbf{运动异常诊断：}

\begin{verbatim}
运动异常
    |
    v
  机器人摔倒了?
    |
  是/        \否
  v            v
 检查电压    运动不流畅?
 get_voltage()   |
    |         是/     \否
    v         v         v
 < 42V?    检查温度    指令无效?
    |      temperature    |
  是/ \否     |          v
  v     v    > 70C?   检查 Arbiter
 低压   检查   |      (遥控器是否在线)
 保护  关节限位
\end{verbatim}

\section*{附录：安全相关环境变量}

\begin{longtable}{lll}
\toprule
环境变量 & 默认值 & 说明 \\
\midrule
\endhead
\texttt{HAN\_DOG\_JOINT\_LIMIT\_RAD} & 3.14 & 关节绝对限位 (rad) \\
\texttt{HAN\_DOG\_ARBITER\_TIMEOUT} & 3 & 遥控器释放后 gRPC 获取控制权等待时间 (s) \\
\texttt{HAN\_DOG\_SHUTDOWN\_TIMEOUT} & 8 & 关机等待超时 (s) \\
\texttt{HAN\_DOG\_STARTUP\_TIMEOUT} & 10 & FSM 启动等待 Grounded 超时 (s) \\
\texttt{HAN\_DOG\_SENSOR\_LOW\_THRESHOLD} & 3 & 传感器低阈值（连续低读数计数） \\
\bottomrule
\end{longtable}

\section*{附录：安全操作总结}

\begin{verbatim}
+------------------------------------------------------+
|              Thunder 安全操作五条铁律                   |
|                                                      |
|  1. 永远带遥控器 -- 遥控器是最终安全手段               |
|  2. 先坐再断电   -- disable() 前必须 sit_down()       |
|  3. 关注电压     -- 低于 44V 准备收工                  |
|  4. 监控温度     -- 超过 70C 立即休息                  |
|  5. 地面运行     -- 永远在地面操作，永远不上桌          |
+------------------------------------------------------+
\end{verbatim}



% ============================================================
%  第 8 章：常见问题
% ============================================================
\chapter{常见问题}
\label{ch:faq}

本章汇总使用 Thunder Python SDK 过程中的常见问题及解决方案。

\section{连接问题}

\subsection{Q: 连接超时 (DEADLINE\_EXCEEDED)}

\noindent\textbf{现象}：调用任何 SDK 方法时抛出 \texttt{grpc.RpcError}，状态码 \texttt{DEADLINE\_EXCEEDED}。

\noindent\textbf{排查步骤}：

\begin{lstlisting}[style=bash]
# 1. 确认网络可达
ping 192.168.66.190

# 2. 确认 gRPC 端口开放
# Linux/macOS:
nc -zv 192.168.66.190 13145
# Windows:
Test-NetConnection -ComputerName 192.168.66.190 -Port 13145

# 3. 确认 brainstem 服务正在运行
ssh sunrise@192.168.66.190 "systemctl status brainstem"
\end{lstlisting}

\noindent\textbf{常见原因与解决}：

\begin{longtable}{ll}
\toprule
原因 & 解决 \\
\midrule
\endhead
机器人未开机 & 打开电池开关，等待系统启动（约 30 秒） \\
开发机和机器人不在同一网段 & 确认两台设备在同一局域网 \\
brainstem 服务未启动 & SSH 登录后运行 \texttt{systemctl start brainstem} \\
防火墙拦截 & 确认开发机防火墙放行 13145 端口 \\
超时时间太短 & 构造时增大 timeout：\texttt{ThunderClient(host, timeout=10.0)} \\
\bottomrule
\end{longtable}

\subsection{Q: 连接被拒绝 (UNAVAILABLE)}

\noindent\textbf{现象}：\texttt{grpc.RpcError} 状态码 \texttt{UNAVAILABLE}，提示 \texttt{Connection refused}。

\noindent\textbf{排查}：

\begin{lstlisting}[style=python]
import grpc

try:
    dog = ThunderClient("192.168.66.190")
    dog.get_state()
except grpc.RpcError as e:
    print(f"状态码: {e.code()}")
    print(f"详情: {e.details()}")
\end{lstlisting}

\noindent\textbf{常见原因}：
\begin{itemize}[nosep]
  \item brainstem 服务未启动或崩溃 -- 检查 \texttt{systemctl status brainstem}
  \item 端口被其他服务占用 -- 检查 \texttt{ss -tlnp | grep 13145}
  \item IP 地址错误 -- 确认机器人实际 IP
\end{itemize}

\subsection{Q: IP 地址不确定}

如果不确定机器人 IP：

\begin{lstlisting}[style=bash]
# 方法 1: 通过路由器管理页面查看 DHCP 租约

# 方法 2: 在已知网段扫描 gRPC 端口
# Linux (需安装 nmap):
nmap -p 13145 192.168.66.0/24

# 方法 3: 如果有 USB 串口连接
# 登录后查看 IP:
ip addr show
\end{lstlisting}

默认 IP 为 \texttt{192.168.66.190}，如果网络配置未更改，直接使用即可。

\section{电机不动}

\subsection{Q: 调用了 enable() 但电机没反应}

\noindent\textbf{排查清单}：

\begin{lstlisting}[style=python]
# 1. 确认状态
state = dog.get_state()
print(f"状态: {state}")
# 应该是 "grounded" 或 "zero"

# 2. 检查电机在线状态
motors = dog.get_motor_status()
offline = [m for m in motors if not m.online]
if offline:
    print(f"离线电机: {[m.id for m in offline]}")

# 3. 检查故障
faulted = [m for m in motors if m.errors]
if faulted:
    print(f"故障电机: {[(m.id, m.errors) for m in faulted]}")
    # 尝试清除
    dog.clear_motor_fault()

# 4. 检查电压
voltages = dog.get_voltage()
print(f"电压范围: {min(voltages):.1f}V ~ {max(voltages):.1f}V")
\end{lstlisting}

\noindent\textbf{常见原因}：

\begin{longtable}{ll}
\toprule
原因 & 解决 \\
\midrule
\endhead
未调用 \texttt{enable()} & 先调用 \texttt{dog.enable()} \\
电机有故障未清除 & \texttt{dog.clear\_motor\_fault()} \\
CAN 线缆松动 & 检查并重新插紧 CAN 接头 \\
电池电量不足 & 充电至 44V 以上 \\
遥控器占用控制权 & 松开遥控器摇杆，等待 3 秒 \\
\bottomrule
\end{longtable}

\subsection{Q: 只有部分电机动，其他不动}

\noindent\textbf{排查}：

\begin{lstlisting}[style=python]
motors = dog.get_motor_status()
for m in motors:
    status = "OK" if m.online and not m.errors else "FAIL"
    print(f"Joint {m.id:2d}: {status}  online={m.online}  errors={m.errors}")
\end{lstlisting}

\noindent\textbf{常见原因}：
\begin{itemize}[nosep]
  \item 对应腿的 CAN 线缆断开 -- 检查该腿的 CAN 总线物理连接
  \item 单个电机驱动板故障 -- 更换后重试
  \item 电机过热保护 -- 查看温度值，等待冷却
\end{itemize}

\subsection{Q: enable() 后状态没变化}

\texttt{enable()} 使能电机，但不会改变 FSM 状态。使能后需要手动切换状态：

\begin{lstlisting}[style=python]
dog.enable()
# 此时状态仍为 "grounded"，电机通电但保持坐姿
dog.stand_up()
# 现在状态变为 "standing"
\end{lstlisting}

\section{机器人摔倒}

\subsection{Q: 行走中突然坐下}

\noindent\textbf{最可能原因}：低压保护触发。

\begin{lstlisting}[style=python]
# 检查电压历史
voltages = dog.get_voltage()
print(f"当前电压: {min(voltages):.1f}V")
# 如果 < 42V，说明是低压保护
\end{lstlisting}

\noindent\textbf{其他可能原因}：

\begin{longtable}{ll}
\toprule
原因 & 诊断方法 \\
\midrule
\endhead
低电压 ($<$ 42V) & \texttt{get\_voltage()} 返回值低于 42 \\
电机过热 & \texttt{get\_motor\_status()} 中 temperature $>$ 80\textdegree C \\
关节超限 & \texttt{get\_motor\_status()} 中 position 绝对值过大 \\
CAN 通信中断 & \texttt{get\_motor\_status()} 中部分电机 online=False \\
传感器断开 & 查看 brainstem 日志中的 SEVERE 信息 \\
\bottomrule
\end{longtable}

\noindent\textbf{查看系统日志}：

\begin{lstlisting}[style=bash]
# SSH 登录机器人查看日志
ssh sunrise@192.168.66.190
# 日志文件在 logs/ 目录下
ls logs/
cat logs/han_dog_$(date +%Y%m%d).log | tail -50
\end{lstlisting}

\subsection{Q: 站起来后立即摔倒}

\noindent\textbf{排查}：
\begin{enumerate}[nosep]
  \item 地面是否太滑 -- 在有摩擦力的地面（橡胶垫、地毯）上测试
  \item 零位是否正确 -- 如果关节零位偏移，站立姿态会畸形
  \item 策略是否匹配 -- 确认加载的运动策略与当前机器人硬件配置一致
\end{enumerate}

\section{行走指令无效}

\subsection{Q: walk() 调用成功但机器人不动}

\noindent\textbf{排查顺序}：

\begin{lstlisting}[style=python]
# 1. 检查状态 -- 必须是 Standing 或 Walking
state = dog.get_state()
print(f"状态: {state}")  # 必须是 "standing" 或 "walking"

# 2. 如果状态不对
if state == "grounded":
    print("需要先站起来: dog.stand_up()")
elif state == "zero":
    print("需要先使能: dog.enable() -> dog.stand_up()")
elif state == "transitioning":
    print("正在过渡，请等待...")
\end{lstlisting}

\noindent\textbf{常见原因}：

\begin{longtable}{ll}
\toprule
原因 & 解决 \\
\midrule
\endhead
不在 Standing 状态 & 先调用 \texttt{stand\_up()} \\
遥控器占用控制权 & 松开遥控器摇杆，等待 3 秒超时 \\
速度参数为 0 & 检查参数：\texttt{dog.walk(vx=0.3)} \\
策略不支持该运动 & 检查当前策略是否支持横移或旋转 \\
\bottomrule
\end{longtable}

\subsection{Q: 遥控器在线时 SDK 指令被忽略}

这是\textbf{正常行为}。ControlArbiter 保证遥控器始终具有最高优先级。

\noindent\textbf{解决方案}：
\begin{enumerate}[nosep]
  \item 关闭遥控器
  \item 或松开遥控器摇杆，等待 3 秒后 SDK 自动获得控制权
\end{enumerate}

\begin{lstlisting}[style=python]
import time

# 等待控制权释放
print("等待遥控器释放控制权...")
time.sleep(4)  # 等待超过 arbiter timeout (默认 3 秒)
dog.walk(vx=0.3)  # 现在应该生效
\end{lstlisting}

\section{策略切换失败}

\subsection{Q: switch\_profile() 报错}

\noindent\textbf{前置条件}：策略切换只能在 \textbf{Grounded} 状态下执行。

\begin{lstlisting}[style=python]
state = dog.get_state()
if state != "grounded":
    print(f"当前状态 {state}，需要先坐下。")
    dog.walk()       # 停步
    dog.sit_down()   # 坐下
    # 等待到 grounded
    import time
    for _ in range(50):
        if dog.get_state() == "grounded":
            break
        time.sleep(0.2)

# 现在可以切换
profile = dog.switch_profile("thunder2v1")
print(f"已切换到: {profile.current}")
\end{lstlisting}

\subsection{Q: 策略名称不存在}

\begin{lstlisting}[style=python]
# 先查看可用策略
profile = dog.get_profile()
print(f"可用策略: {profile.available}")
print(f"策略说明: {profile.descriptions}")

# 使用列表中的名称
dog.switch_profile(profile.available[0])
\end{lstlisting}

\subsection{Q: 切换策略后行为异常}

每个策略针对特定硬件配置和运动模式进行训练。切换策略后如果出现异常运动：
\begin{enumerate}[nosep]
  \item 确认策略与当前硬件配置匹配
  \item 立即坐下：\texttt{dog.sit\_down()}
  \item 切回之前正常工作的策略
\end{enumerate}

\section{IMU 数据异常}

\subsection{Q: 四元数数值看起来不对}

\noindent\textbf{Hamilton vs ROS 约定}：

Thunder SDK 使用 Hamilton 约定：\texttt{(w, x, y, z)}

\begin{lstlisting}[style=python]
imu = next(iter(dog.listen_imu()))
# SDK 返回: (w, x, y, z)
print(f"w={imu.quaternion.w}, x={imu.quaternion.x}, "
      f"y={imu.quaternion.y}, z={imu.quaternion.z}")
\end{lstlisting}

如果你的应用使用 ROS 约定 \texttt{(x, y, z, w)}，需要手动转换：

\begin{lstlisting}[style=python]
# SDK (Hamilton) -> ROS
ros_quat = (imu.quaternion.x, imu.quaternion.y,
            imu.quaternion.z, imu.quaternion.w)

# 如果使用 scipy
from scipy.spatial.transform import Rotation
r = Rotation.from_quat([
    imu.quaternion.x,
    imu.quaternion.y,
    imu.quaternion.z,
    imu.quaternion.w,
])  # scipy 使用 (x,y,z,w) 顺序
euler = r.as_euler('xyz', degrees=True)
\end{lstlisting}

\subsection{Q: 机器人静止但角速度不为零}

IMU 的陀螺仪存在零偏（bias），静止时读数不会严格为零。典型零偏范围：

\begin{verbatim}
|gyro.x| < 0.02 rad/s  -- 正常
|gyro.x| > 0.05 rad/s  -- 可能需要标定
\end{verbatim}

如需高精度角速度，可在启动时采集 100 帧静止数据取均值作为零偏，后续数据减去此偏移。

\subsection{Q: listen\_imu() 频率不是 50Hz}

IMU 标称频率约 50Hz，实际可能在 48--52Hz 波动。这是正常的。如果频率显著偏低（$<$ 30Hz），检查：
\begin{itemize}[nosep]
  \item 串口连接是否正常
  \item brainstem 服务是否负载过高
  \item 网络延迟（gRPC 传输开销）
\end{itemize}

\section{性能与延迟}

\subsection{Q: gRPC 调用延迟有多大？}

\begin{longtable}{lll}
\toprule
调用类型 & 典型延迟 & 说明 \\
\midrule
\endhead
单次查询 (\texttt{get\_state()} 等) & 1--5 ms & 局域网内 \\
运动指令 (\texttt{walk()} 等) & 1--5 ms & 局域网内 \\
首次连接 & 50--200 ms & 包含 HTTP/2 握手 \\
流式数据 (\texttt{listen\_imu()} 等) & $<$ 1 ms/帧 & 建立后持续推送 \\
\bottomrule
\end{longtable}

\subsection{Q: 50Hz 控制频率够用吗？}

Thunder 的运动控制循环固定为 50Hz（20ms 周期）。SDK 发送的指令在下一个控制周期生效。这意味着：
\begin{itemize}[nosep]
  \item 指令最大延迟：20ms（一个控制周期）
  \item 不需要以超过 50Hz 的频率发送 walk() 指令
  \item 建议的 walk() 指令发送频率：10--50Hz
\end{itemize}

\begin{lstlisting}[style=python]
import time

# 不需要这样做（浪费资源）:
while True:
    dog.walk(vx=0.3)
    time.sleep(0.001)  # 1000Hz -- 过高，无意义

# 推荐做法:
while True:
    dog.walk(vx=0.3)
    time.sleep(0.05)  # 20Hz -- 足够
\end{lstlisting}

\subsection{Q: Python GIL 会影响实时性吗？}

Python 的全局解释器锁（GIL）对 gRPC 调用影响很小，因为 gRPC 底层使用 C 扩展进行网络 I/O，不受 GIL 限制。但以下场景需要注意：

\begin{itemize}[nosep]
  \item \textbf{多线程同时调用}：gRPC 客户端本身是线程安全的，但应避免多线程同时发送运动指令（语义冲突）
  \item \textbf{数据处理密集}：如果在 \texttt{listen\_imu()} 回调中做大量 numpy 计算，可能导致帧处理延迟。建议将数据处理放到独立线程
\end{itemize}

\begin{lstlisting}[style=python]
import threading
import queue

# 推荐模式：生产者-消费者
data_queue = queue.Queue(maxsize=100)

def imu_reader(dog):
    """生产者：读取 IMU 数据放入队列。"""
    for imu in dog.listen_imu():
        try:
            data_queue.put_nowait(imu)
        except queue.Full:
            data_queue.get()  # 丢弃最旧的
            data_queue.put_nowait(imu)

def data_processor():
    """消费者：处理 IMU 数据。"""
    while True:
        imu = data_queue.get()
        # 在这里做复杂计算...
        pass

reader = threading.Thread(target=imu_reader, args=(dog,), daemon=True)
processor = threading.Thread(target=data_processor, daemon=True)
reader.start()
processor.start()
\end{lstlisting}

\section{环境变量参考}

以下环境变量在机器人端（brainstem 服务）生效，SDK 客户端不直接使用。了解这些变量有助于理解系统行为和排障。

{\small
\begin{longtable}{lllll}
\toprule
环境变量 & 类型 & 默认值 & 有效范围 & 说明 \\
\midrule
\endhead
\texttt{HAN\_DOG\_PORT} & int & 13145 & 1--65535 & gRPC 服务监听端口 \\
\texttt{HAN\_DOG\_IMU\_PORT} & str & /dev/ttyUSB1 & 设备路径 & IMU 串口设备 \\
\texttt{HAN\_DOG\_YUNZHUO\_PORT} & str & /dev/yunzhuo & 设备路径 & YUNZHUO 遥控器串口 \\
\texttt{HAN\_DOG\_XBOX\_DEVICE} & str & /dev/input/js0 & 设备路径 & Xbox 手柄设备 \\
\texttt{HAN\_DOG\_XBOX\_CONFIG} & str & config/xbox.json & 文件路径 & Xbox 手柄按键映射 \\
\texttt{HAN\_DOG\_DEFAULT\_PROFILE} & str & (无) & 策略名称 & 默认加载的运动策略 \\
\texttt{HAN\_DOG\_PROFILE\_DIR} & str & profiles/ & 目录路径 & 策略文件目录 \\
\texttt{HAN\_DOG\_ARBITER\_TIMEOUT} & int & 3 & 1--60 & 遥控器释放后控制权切换 (s) \\
\texttt{HAN\_DOG\_SENSOR\_LOW\_THRESHOLD} & int & 3 & $\geq$ 1 & 传感器低读数计数阈值 \\
\texttt{HAN\_DOG\_SHUTDOWN\_TIMEOUT} & int & 8 & 1--300 & 关机等待超时 (s) \\
\texttt{HAN\_DOG\_STARTUP\_TIMEOUT} & int & 10 & 1--300 & 启动等待 Grounded 超时 (s) \\
\texttt{HAN\_DOG\_JOINT\_LIMIT\_RAD} & float & 3.14 & $>$ 0 & 关节绝对限位 (rad) \\
\texttt{HAN\_DOG\_LOG\_DIR} & str & logs & 目录路径 & 日志文件目录 \\
\texttt{HAN\_DOG\_LOG} & str & INFO & Logger 级别 & 日志级别 \\
\texttt{HAN\_DOG\_DEBUG\_TUI} & str & (无) & true/其他 & 启用调试 TUI \\
\texttt{MEDULLA\_STANDALONE} & str & (无) & 1/其他 & 启用独立模式 \\
\bottomrule
\end{longtable}
}

\section{网络排障}

\subsection{基本连通性检查}

\begin{lstlisting}[style=bash]
# 第 1 步: Ping 机器人
ping 192.168.66.190
# 预期: 响应时间 < 5ms (局域网)

# 第 2 步: 检查 gRPC 端口
# Linux/macOS:
nc -zv 192.168.66.190 13145
# 预期输出: Connection to 192.168.66.190 13145 port [tcp/*] succeeded!

# Windows PowerShell:
Test-NetConnection -ComputerName 192.168.66.190 -Port 13145
# 预期: TcpTestSucceeded : True

# 第 3 步: 检查服务状态 (SSH)
ssh sunrise@192.168.66.190 "systemctl status brainstem"
# 预期: Active: active (running)
\end{lstlisting}

\subsection{防火墙检查}

\begin{lstlisting}[style=bash]
# Linux 开发机
sudo iptables -L -n | grep 13145

# Windows 开发机
netsh advfirewall firewall show rule name=all | findstr 13145

# 临时放行 (Linux)
sudo iptables -A OUTPUT -p tcp --dport 13145 -j ACCEPT

# 临时放行 (Windows PowerShell, 管理员)
New-NetFirewallRule -DisplayName "Thunder gRPC" -Direction Outbound -Protocol TCP -RemotePort 13145 -Action Allow
\end{lstlisting}

\subsection{网络拓扑要求}

\begin{verbatim}
开发机                    Thunder 机器人
(192.168.66.x)           (192.168.66.190)
     |                         |
     +--- 同一局域网/交换机 ---+
           (无 NAT、无代理)
\end{verbatim}

\begin{itemize}[nosep]
  \item SDK 使用 gRPC over HTTP/2，需要直连（不经过 HTTP 代理）
  \item 不支持跨 NAT 连接（如果开发机在不同子网，需要配置路由）
  \item Wi-Fi 连接可能有延迟抖动，建议使用有线网络
\end{itemize}

\subsection{远程开发（通过 SSH 隧道）}

如果开发机不在局域网内，可通过 SSH 隧道连接：

\begin{lstlisting}[style=bash]
# 在有局域网访问的跳板机上建立隧道
ssh -L 13145:192.168.66.190:13145 sunrise@192.168.66.190

# 然后 SDK 连接 localhost
dog = ThunderClient("127.0.0.1", port=13145)
\end{lstlisting}

\section{提交 Bug 报告}

遇到无法自行解决的问题时，请按以下格式提交 Bug 报告。

\subsection{必要信息}

\begin{lstlisting}[style=plain,numbers=none]
1. SDK 版本:
   python -c "import brainstem_sdk; print(brainstem_sdk.__version__)"

2. Python 版本:
   python --version

3. 操作系统:
   (例: Ubuntu 22.04 / Windows 11 / macOS 14)

4. 机器人固件版本:
   (SSH 登录后查看 brainstem 版本)

5. 问题描述:
   (清晰描述预期行为和实际行为)

6. 复现步骤:
   (从连接开始，逐步列出操作)

7. 错误信息:
   (完整的异常堆栈或错误输出)

8. 机器人日志:
   (SSH 登录，提供 logs/han_dog_YYYYMMDD.log 中相关片段)
\end{lstlisting}

\subsection{收集诊断信息的脚本}

\begin{lstlisting}[style=python]
"""诊断信息收集脚本。将输出复制到 Bug 报告中。"""

import sys
import platform
import grpc

try:
    import brainstem_sdk
    sdk_ver = brainstem_sdk.__version__
except Exception as e:
    sdk_ver = f"导入失败: {e}"

print("=" * 50)
print("Thunder SDK 诊断信息")
print("=" * 50)
print(f"SDK 版本:    {sdk_ver}")
print(f"Python:      {sys.version}")
print(f"平台:        {platform.platform()}")
print(f"grpc 版本:   {grpc.__version__}")
print()

host = input("机器人 IP (默认 192.168.66.190): ").strip() or "192.168.66.190"

try:
    from brainstem_sdk import ThunderClient
    dog = ThunderClient(host, timeout=5.0)

    # 状态
    state = dog.get_state()
    print(f"连接:        成功")
    print(f"机器人状态:  {state}")

    # 电压
    voltages = dog.get_voltage()
    print(f"电压范围:    {min(voltages):.1f}V ~ {max(voltages):.1f}V")

    # 电机
    motors = dog.get_motor_status()
    online_count = sum(1 for m in motors if m.online)
    fault_count = sum(1 for m in motors if m.errors)
    max_temp = max(m.temperature for m in motors)
    print(f"电机在线:    {online_count}/16")
    print(f"电机故障:    {fault_count}/16")
    print(f"最高温度:    {max_temp:.0f}C")

    # 故障详情
    for m in motors:
        if m.errors or not m.online:
            print(f"  Joint {m.id}: online={m.online} errors={m.errors} "
                  f"temp={m.temperature:.0f}C voltage={m.voltage:.1f}V")

    # 策略
    profile = dog.get_profile()
    print(f"当前策略:    {profile.current}")
    print(f"可用策略:    {profile.available}")

    dog.close()

except grpc.RpcError as e:
    print(f"连接:        失败")
    print(f"错误码:      {e.code()}")
    print(f"详情:        {e.details()}")
except Exception as e:
    print(f"异常:        {type(e).__name__}: {e}")

print()
print("=" * 50)
print("请将以上输出复制到 Bug 报告中。")
\end{lstlisting}

\subsection{日志收集}

\begin{lstlisting}[style=bash]
# SSH 登录机器人
ssh sunrise@192.168.66.190

# 查看今天的日志
cat logs/han_dog_$(date +%Y%m%d).log

# 只看错误和警告
grep -E "WARNING|SEVERE" logs/han_dog_$(date +%Y%m%d).log

# 查看最近 100 行
tail -100 logs/han_dog_$(date +%Y%m%d).log

# 导出日志到本地（在开发机上执行）
scp sunrise@192.168.66.190:logs/han_dog_*.log ./thunder_logs/
\end{lstlisting}

\subsection{联系方式}

\begin{itemize}[nosep]
  \item GitHub Issues：在 brainstem 仓库提交 Issue
  \item 附上诊断脚本输出和相关日志
\end{itemize}

\section*{附录：快速排障速查表}

{\small
\begin{longtable}{llll}
\toprule
症状 & 首先检查 & 可能原因 & 解决方案 \\
\midrule
\endhead
连接超时 & \texttt{ping} 机器人 & 网络不通 & 检查网线/Wi-Fi/IP \\
连接被拒 & brainstem 服务状态 & 服务未启动 & \texttt{systemctl start brainstem} \\
enable() 无效 & \texttt{get\_motor\_status()} & 电机离线/有故障 & 检查 CAN 线，清除故障 \\
walk() 无效 & \texttt{get\_state()} & 不在 Standing & 先 \texttt{stand\_up()} \\
walk() 被忽略 & 遥控器状态 & Arbiter 被占用 & 松开遥控器，等 3 秒 \\
突然坐下 & \texttt{get\_voltage()} & 低压保护 ($<$ 42V) & 充电 \\
突然摔倒 & 电机温度 & 过热/关节超限 & 冷却 / 重新标零 \\
策略切换失败 & \texttt{get\_state()} & 不在 Grounded & 先 \texttt{sit\_down()} \\
IMU 数据异常 & 四元数约定 & Hamilton vs ROS & 调整字段顺序 \\
数据流中断 & 网络连接 & gRPC 通道断开 & 重建 ThunderClient \\
多客户端冲突 & 客户端数量 & 多 SDK 实例 & 保证单客户端 \\
\bottomrule
\end{longtable}
}

% ============================================================
%  文档结束
% ============================================================

\end{document}
